\documentclass[11pt,oneside]{article}
\usepackage[a4paper,margin=25mm,headheight=15pt]{geometry}
\usepackage[T1]{fontenc}
\usepackage[utf8]{inputenc}
\usepackage{lmodern}
\usepackage{microtype}
\usepackage{amsmath,amssymb,amsthm,mathtools}
\usepackage{booktabs,longtable,tabularx,array}
\usepackage{enumitem}
\usepackage{xcolor}
\usepackage{listings}
\usepackage{hyperref}
\usepackage[nameinlink,noabbrev]{cleveref}
\usepackage{fancyhdr}
\usepackage{lastpage}
\usepackage{url}
\usepackage{graphicx}
\usepackage{etoolbox}

\definecolor{scopeblue}{RGB}{31,78,121}
\definecolor{scopegray}{RGB}{245,247,249}
\definecolor{scopegreen}{RGB}{34,110,74}
\definecolor{scopered}{RGB}{140,45,45}
\hypersetup{
  colorlinks=true,
  linkcolor=scopeblue,
  citecolor=scopegreen,
  urlcolor=scopeblue,
  pdftitle={SCOPE-iPRA: Scoped Certificates, Observational Projection, and Explicit Residuals},
  pdfauthor={OpenAI, prepared for Vladimir Reshetnikov}
}

\pagestyle{fancy}
\fancyhf{}
\fancyhead[L]{\small SCOPE--iPRA}
\fancyhead[R]{\small 16 November 2026}
\fancyfoot[C]{\small \thepage\ of \pageref*{LastPage}}

\setlist{nosep,leftmargin=2em}
\setlength{\parindent}{0pt}
\setlength{\parskip}{0.55em}
\setcounter{tocdepth}{2}
\setcounter{secnumdepth}{3}

\newtheorem{theorem}{Theorem}[section]
\newtheorem{lemma}[theorem]{Lemma}
\newtheorem{proposition}[theorem]{Proposition}
\newtheorem{corollary}[theorem]{Corollary}
\newtheorem{definition}[theorem]{Definition}
\newtheorem{principle}[theorem]{Principle}
\newtheorem{warning}[theorem]{Warning}
\theoremstyle{remark}
\newtheorem{remark}[theorem]{Remark}
\newtheorem{example}[theorem]{Example}

\newcommand{\ipra}{\mathsf{iPRA}}
\newcommand{\PRAAx}{\mathsf{PRAAx}}
\newcommand{\Prf}{\mathsf{Prf}}
\newcommand{\SCOPE}{\mathsf{SCOPE}}
\newcommand{\PTm}{\mathsf{PTm}}
\newcommand{\Fm}{\mathsf{Fm}}
\newcommand{\PR}{\mathsf{PR}}
\newcommand{\Valid}{\mathsf{Valid}}
\newcommand{\Obs}{\mathsf{Obs}}
\newcommand{\Goal}{\mathsf{Goal}}
\newcommand{\Res}{\mathsf{Res}}
\newcommand{\bit}{\mathsf{bit}}
\newcommand{\rank}{\mathsf{rank}}
\newcommand{\Nat}{\mathbb N}
\newcommand{\code}[1]{\texttt{\detokenize{#1}}}
\newcommand{\seq}[3][]{#2\;\vdash_{\PRAAx}^{#1}\;#3}
\newcommand{\checked}{\textcolor{scopegreen}{\textbf{inherited checked}}}
\newcommand{\human}{\textcolor{scopeblue}{\textbf{new human proof}}}
\newcommand{\blueprint}{\textcolor{scopered}{\textbf{Lean blueprint}}}

\lstdefinelanguage{Lean}{
  morekeywords={theorem,def,structure,inductive,namespace,end,variable,open,import,where,by,exact,apply,have,let,match,with,if,then,else,fun,forall,exists},
  sensitive=true,
  morecomment=[l]{--},
  morecomment=[s]{/-}{-/},
  morestring=[b]"
}
\lstset{
  language=Lean,
  basicstyle=\ttfamily\small,
  keywordstyle=\color{scopeblue}\bfseries,
  commentstyle=\color{scopegreen},
  stringstyle=\color{scopered},
  backgroundcolor=\color{scopegray},
  frame=single,
  framerule=0.2pt,
  breaklines=true,
  columns=fullflexible,
  keepspaces=true,
  showstringspaces=false
}

\title{\textbf{SCOPE--iPRA}\\[2mm]
\Large Scoped Certificates, Observational Projection, and Explicit Residuals\\[2mm]
\large A general framework for embedding mathematics and proofs into intuitionistic primitive recursive arithmetic}
\author{Prepared for Vladimir Reshetnikov\\OpenAI}
\date{16 November 2026\\Repository baseline: \code{58627ed89814ba122e1f0f012f1aba103cda267b}}

\begin{document}
\maketitle

\begin{abstract}
The ProveIt development now contains a kernel-checked ordinary presentation of intuitionistic primitive recursive arithmetic, a substantial library of checked object-theory derivations, and more than one hundred research packets that translate classical mathematics into finite quantitative claims.  The remaining general difficulty is not a lack of encodings.  It is the absence of one discipline that says exactly when a finite witness may escape its proof scope, how cyclic precision requirements are closed before induction, how bounded quantifier alternation is reflected without quantified induction predicates, and how unresolved mathematical producers remain visible in the final sequent.

This report develops \emph{SCOPE--iPRA}: Scoped Certificates, Observational Projection, and Explicit Residuals.  Every fixed mathematical argument is compiled at a finite query.  Fixed bounded formulas are translated to primitive-recursive characteristic terms; finite loops are reflected by quantifier-free induction; cyclic demand networks are closed by an explicit least-fixed-point iteration on a capped product; non-unique witnesses are handled by backward viability; and proof-local certificates are eliminated by ordinary existential elimination.  A certificate may become public only as an existential output, through a certificate-independent exact/quotient/metric observation, by least bounded canonical naming, or as an explicit residual producer obligation.  These five escape modes yield a scope-normal-form theorem whose endpoint is a literal object derivation
\[
  B_1,\ldots,B_r,\Gamma\ \vdash_{\PRAAx}\ \Phi.
\]
Supplying derivations of the residuals closes the theorem by the already checked simultaneous-cut rule.

The report gives complete human proofs of bounded reflection, guarded least choice, fixed-trace reflection, capped cyclic demand closure, finite strategy extraction, backward viability, observer elimination, finite observational descent, height-bound/list equivalence, and the SCOPE normalization theorem.  It maps each construction to the current \code{ConciseIPRA}, \code{IPRADerivations}, and \code{IPRABridge} APIs.  The archive also contains a no-placeholder structural Lean candidate, detailed Lean blueprints, exact Python reference algorithms, fourteen passing tests, a worked certificate, a theorem ledger, and a difficulty matrix.  The new generic theorems are not claimed kernel-checked in this archive.
\end{abstract}

\tableofcontents
\newpage

\section{Executive summary and exact status}

\subsection{The repository situation}

The inspected ProveIt snapshot separates three kinds of material that must not be conflated.

\begin{enumerate}
\item \textbf{Checked logical substrate.}  \code{ConciseIPRA} defines scoped terms, formulas, intuitionistic natural deduction with equality, primitive-recursive function symbols, the ordinary axiom family \(\PRAAx\), and checked translations to smaller native calculi.
\item \textbf{Checked object-theory mathematics.}  \code{IPRADerivations} contains Lean-kernel-checked values of types such as
\[
  ([] : \mathsf{List}(\Fm\,\PTm\,0))\ \vdash_{\PRAAx}\ \varphi,
\]
and open variants under arbitrary finite contexts.  Its reusable core already includes arithmetic, order, pairing, bounded search, division, finite sums and products, descent, and finite-dimensional linear algebra.
\item \textbf{Research packets.}  \code{docs/proofs} contains a much larger set of delivered human proofs, executable finite checkers, and Lean plans.  The catalog explicitly records that these packets were not collectively elaborated or verified.  Their conceptual content is valuable, but a prose theorem or a semantic Lean proposition is not silently promoted to a \(\Prf\,\PRAAx\) derivation.
\end{enumerate}

The foundations packets already identify most of the right ingredients: finite proof graphs, bounded reflection, least search, reverse demand propagation, trace reflection, relational witness clouds, stable observations, explicit residuals, and scoped certificate elimination.  Their overlap is substantial.  The current problem is to consolidate them into one normal form aligned with the checked library, remove obsolete open tasks, and make the remaining producer boundaries mathematically exact.

\subsection{The SCOPE contribution}

SCOPE stands for \emph{Scoped Certificates, Observational Projection, and Explicit Residuals}.  It adds one organizing principle to the earlier finite-contract frameworks:

\begin{quote}
A certificate is a local proof resource.  It is not public mathematical data unless the final theorem explicitly chooses one of five justified escape modes.
\end{quote}

The five modes are:

\begin{enumerate}
\item \emph{private}: eliminate the certificate by existential elimination;
\item \emph{existential}: use certificate-dependent data only as the witness of an outer existential;
\item \emph{stable}: export an exact, quotient, metric, or gap-stabilized observation proved independent of the certificate in the required sense;
\item \emph{canonical}: name the least accepted code by checked bounded search;
\item \emph{residual}: if no producer is available, retain its exact existence or quantitative data formula as a hypothesis.
\end{enumerate}

This classification is not cosmetic.  It separates four notions that are repeatedly confused in constructive formalization: existence of a witness, the right to name a witness, invariance of a finite observation, and availability of a theorem-specific producer.

\subsection{Principal theorems}

The report proves the following generic results at the human mathematical level.

\begin{description}[style=nextline]
\item[Bounded reflection.] Every fixed bounded arithmetic formula built from primitive-recursive terms has a primitive-recursive characteristic bit.  iPRA proves the formula equivalent to the bit being one, and proves its decidability.  Nested bounded quantifiers are handled by finite folds and least counterexample/witness search; induction is only on quantifier-free equations.

\item[Guarded canonical choice.] A bounded existential has a primitive-recursive least witness together with a success bit and the bound as failure sentinel.  Naming is therefore available exactly when a finite bound is supplied.

\item[Fixed-trace reflection.] A fixed primitive-recursive state machine with a quantifier-free invariant, object-theory base proof, step proof, and terminal extractor yields an object-theory theorem at arbitrary primitive-recursive fuel.

\item[Capped cyclic demand closure.] For a monotone demand transformer on a finite product \(\prod_{i<k}[0,C_i]\), iteration of \(d\mapsto d\vee F(d)\) from zero reaches the least prefixed point after at most \(\sum_i C_i\) strict increases.  A rank proof is directly arithmetizable by quantifier-free induction.

\item[Finite strategy extraction.] Every fixed bounded quantifier game admits a finite strategy verifier and a dual counterstrategy verifier.  Backward evaluation constructs one of them constructively; no appeal to infinite game determinacy is involved.

\item[Backward viability.] For finite levels of non-unique approximants, a compatible path exists exactly when the initial backward-viability bit is one.  The lexicographically least \emph{viable} path is primitive recursive.  Choosing the least local witness before viability is unsound because it may select a dead end.

\item[Observer elimination.] Exact certificate-independent observations descend by equality; quotient observations descend through witnessed congruences; metric observations descend with an explicit diameter; and exact branch labels descend only under an explicit separation gap.

\item[Scope normal form.] Every fixed finite interaction graph satisfying the local contracts compiles to a derivation from precisely its reachable residual producers.  All internal certificates have one of the five escape modes, and no hidden choice remains.
\end{description}

\subsection{What is resolved and what remains conditional}

The framework completely resolves recurring proof-engineering difficulties: context composition, binder-safe specialization, bounded quantifier reflection, least bounded naming, cyclic finite demand closure, finite relational path selection, and certificate scoping.  It also gives principled reformulations of several overstrong requests.  A represented nonzero real is not invertible without apartness; colliding roots do not have continuous ordered labels; metastability is not a Cauchy modulus; and a cardinality bound is not a height cutoff.

The framework deliberately does not manufacture deep mathematical producer data.  General effective height bounds, zero-free regions, exceptional-zero constants, finite covers, branch gaps, or convergence moduli remain theorem-specific residuals until they are proved.  Once supplied, their finite verification and use are generic.

\subsection{Verification labels used in this report}

Four statuses are kept separate throughout:

\begin{center}
\begin{tabularx}{\textwidth}{@{}lX@{}}
\toprule
Label & Meaning \\
\midrule
\checked & The cited declaration is already checked in the inspected ProveIt repository. \\
\human & The present report supplies a complete mathematical argument, but not a new kernel check. \\
\blueprint & The archive supplies definitions and a dependency-ordered Lean target; it was not elaborated here. \\
Executable reference & Python code exactly implements a finite algorithm and its tests pass; it is not an iPRA proof. \\
\bottomrule
\end{tabularx}
\end{center}

The archive makes no claim stronger than these labels permit.

\section{The exact logical substrate}

\subsection{Ordinary iPRA as the canonical interface}

The canonical target is the ordinary deeply embedded judgment
\[
  \Gamma\ \vdash_{\PRAAx}\ \varphi,
\]
implemented by \code{LeanProofs.ConciseIPRA.Prf}.  A sequent lives over a fixed finite number \(n\) of free numerical variables.  Its terms have type \(\PTm\,n\), its formulas have type \(\Fm\,\PTm\,n\), and every formula in \(\Gamma\) has the same scope.

Primitive-recursive function symbols are generated by zero, successor, projections, composition, and primitive recursion.  The term language has one application constructor, so every fixed primitive-recursive algorithm can be represented by a symbol in \(\PR\) and applied to arbitrary object terms.  The axiom family \(\PRAAx\) consists of defining equations for those symbols, successor separation and injectivity, and induction for quantifier-free formulas.

The logical calculus is intuitionistic natural deduction with equality.  The structural operations used by SCOPE are already checked:

\begin{itemize}
\item \code{Prf.subst}: simultaneous substitution of arbitrary object terms for free variables;
\item \code{Prf.weaken}: context weakening;
\item \code{Prf.trans}: simultaneous cut, replacing every assumption by a derivation of it;
\item \code{Prf.exE}: existential elimination with the witness kept under a fresh-variable scope;
\item \code{Prf.eqSubst}: simultaneous Leibniz transport.
\end{itemize}

These operations are sufficient for the entire structural part of SCOPE.  No new logical axiom, choice principle, or reflection rule is needed.

\subsection{Quantifier-free induction and the compilation discipline}

The axiom \code{PRA.qfInd} applies to a formula \(I(t,\vec x)\) only after Lean has established \code{I.QF}.  The checked derivation library exposes a convenient open form \code{LeanProofs.IPRA.ind}.  SCOPE uses this restriction as a design rule rather than treating it as an inconvenience:

\begin{principle}[Equation-first induction]
Any finite property needed inside an induction is compiled to a primitive-recursive bit or state component.  The induction matrix is an equation or a finite quantifier-free combination of equations.  Reflection back to the intended bounded formula happens before or after the induction, never by pretending a bounded quantifier is syntactically quantifier-free.
\end{principle}

For example, ``every table entry below \(t\) is valid'' is represented by the equation \(\mathsf{all}(t,q)=1\).  ``Some transition has failed'' is a failure-accumulator bit.  ``Every coordinate is within its cap'' is either a fixed finite conjunction or a finite-fold bit.  This uniform pattern is what makes trace compilation reusable.

\subsection{Current arithmetic support}

The present \code{IPRADerivations} library is substantially richer than several earlier framework packets assumed.  SCOPE relies on the following checked components.

\begin{itemize}
\item \code{Core.Basic}: arithmetic symbols, defining equations, open induction, and instantiation;
\item \code{Core.Definitions}: \code{PTm.toPR}, its evaluation theorem \code{toPR_eq}, arbitrary-arity \code{recDef}, and finite sums/products;
\item \code{Core.Order}: reflected order, totality, equality decidability, and monotonicity;
\item \code{Core.Pairing}: \code{eqb}, pairing/projections, and simultaneous recursion;
\item \code{Core.Search}: a least-zero search \(\mu\), soundness, foundness, and minimality;
\item \code{Core.Division}, \code{Core.MultiSums}, and \code{Core.Permutation}: the arithmetic needed for finite code spaces;
\item polynomial and integer-difference backends, finite descent, and finite-dimensional linear algebra for application nodes.
\end{itemize}

One reported obstacle is therefore obsolete.  Older bounded-reflection outlines asked for an arity-indexed abstraction turning an open term into a primitive-recursive symbol.  The current declaration
\[
  \code{PTm.toPR} : \PTm\,k \longrightarrow \PR\,k
\]
and its object-theory theorem \code{toPR_eq} provide exactly that operation.  Together with \code{recDef}, they permit arbitrary-parameter Boolean folds without an internal universal evaluator.

\subsection{Fixed specialization, not a universal evaluator}

SCOPE is a schema over a \emph{fixed} finite program or formula.  Lean's metatheory traverses the finite syntax tree, constructs a particular \(\PR\) symbol, and constructs a particular \(\Prf\,\PRAAx\) derivation.  The object theory need not evaluate arbitrary encoded \(\PR\) programs.

This distinction avoids a common impossibility.  A universal total evaluator for all primitive-recursive programs cannot itself be primitive recursive.  Nothing in SCOPE asks for one.  Each theorem has finitely many concrete symbols, and \code{PTm.toPR} compiles those concrete open terms directly.

\subsection{Object proofs versus semantic Lean theorems}

A native Lean theorem about \(\Nat\), lists, rationals, or a recursive function is useful for discovering an invariant and testing an algorithm.  It is not automatically a theorem of the embedded object theory.  Soundness goes from an object derivation to semantics; there is no unrestricted reverse reflection from arbitrary Lean proofs.

Accordingly, every SCOPE node has two logically distinct ingredients:

\begin{enumerate}
\item a host-level finite description used to generate symbols, terms, and scheduling data;
\item an object-level derivation builder whose output type is a \(\Prf\,\PRAAx\) sequent.
\end{enumerate}

The compiler is trusted only as a producer of a proof term that Lean's kernel checks.  External code may propose certificates and can greatly accelerate discovery, but acceptance is ultimately represented by object arithmetic and an object derivation.

\subsection{Backend transport}

The repository's bridge now exposes closed, open, and scoped transports between ordinary ConciseIPRA and Foundation's standard iPRA presentation.  The relevant APIs include \code{standardToConciseSequent}, \code{conciseToStandardSequent}, and scoped round-trip theorems.  Thus open-sequent transport is no longer a missing generic theorem.

SCOPE therefore develops all reusable mathematics once in the ordinary \(\Prf\,\PRAAx\) interface.  Native scalar calculi and Foundation arithmetic are export backends.  This avoids parallel libraries and lets application endpoints choose the presentation they require.

\section{What an embedding of mathematics into iPRA must provide}

\subsection{Finite queries rather than completed infinite objects}

An ordinary mathematical object may be infinite: a real number, a function, a compact space, a probability distribution, a formal power series, or a sheaf.  An iPRA theorem does not need to internalize the completed object if every concrete use asks only a finite question.

A \emph{query} \(q\) records the finite observation requested by a consumer.  Typical components are:

\begin{itemize}
\item a precision \(p\);
\item a finite domain bound \(N\);
\item a requested derivative order or truncation degree;
\item a mesh, cover, or search height;
\item a finite list of evaluation points;
\item a resource budget or proof horizon.
\end{itemize}

The producer responds with finite coded data \(c<C(q)\) and a decidable validity relation \(\Valid(q,c)\).  The theorem may then return a finite observation, a relation of acceptable observations, or a conditional result with explicit producer assumptions.

\subsection{Representations, realizers, and observations}

\begin{definition}[Finitary presentation]
A finitary presentation consists of:
\begin{enumerate}
\item a primitive-recursive query code space;
\item for each query, a bounded certificate space \(c<C(q)\);
\item a fixed bounded or quantifier-free validity formula \(\Valid(q,c)\);
\item a primitive-recursive observation map \(\Obs(q,c)\), possibly relational;
\item object-theory derivations connecting validity to the advertised mathematical specification.
\end{enumerate}
\end{definition}

The presentation need not choose one representative globally.  For a root with colliding branches, \(c\) may encode an unordered interval multiset.  For an eigenspace, it may encode an approximate projector rather than a signed eigenvector.  For a compactness argument, it may encode a finite net and its coverage certificate.  The observation interface should match what is stable under the available data.

\subsection{Residual producers}

Many finite verifiers are elementary once a quantitative producer is supplied.  The difficult mathematical content often lies in proving that the producer exists uniformly.  SCOPE records this distinction in the final sequent.

\begin{definition}[Residual producer]
A residual producer is an object-language formula \(B(q)\) asserting the existence or correctness of quantitative data required by the finite compiler but not constructed by the generic framework.  Examples include:
\begin{itemize}
\item a height cutoff containing all solutions;
\item a finite cover of prescribed mesh;
\item an apartness or branch-separation margin;
\item a convergence or tail modulus;
\item an explicit analytic constant or zero-free region;
\item a rational representative of a descent class.
\end{itemize}
\end{definition}

The generic result has type
\[
  B_1(q),\ldots,B_r(q),\Gamma(q)\ \vdash_{\PRAAx}\ \Goal(q).
\]
If later work proves each \(B_i\), simultaneous cut closes the result.  Until then, the theorem remains a precise conditional theorem rather than an overstated unconditional one.

\subsection{Finite interaction graphs}

A large proof is organized as a fixed finite graph.  Its nodes may be:

\begin{enumerate}
\item exact primitive-recursive transducers;
\item Boolean decision nodes;
\item bounded searches;
\item clocked loops with quantifier-free invariants;
\item finite relational producers;
\item observer or quotient nodes;
\item residual producer leaves.
\end{enumerate}

Edges carry finite data and backward resource demands.  Acyclic data flow is topologically scheduled.  Cyclic precision dependencies are not executed as an unbounded feedback loop; they are first closed in a finite capped demand lattice, as proved in \cref{sec:demand-closure}.  After closure, the forward computation is again a fixed finite trace.

\subsection{The five public endpoint modes}

The same mathematical theorem may have several constructively distinct endpoints.

\paragraph{Exact endpoint.}
Return a concrete term and an exact equation or predicate.  This is appropriate for finite arithmetic, canonical least witnesses, and robustly isolated branches.

\paragraph{Existential endpoint.}
Prove \(\exists z\,P(q,z)\) while allowing the chosen \(z\) to depend on a private certificate.  No global selector is exposed.

\paragraph{Relational endpoint.}
Return a finite cloud \(R(q,z)\) of acceptable answers or a statement that every valid answer has a desired stable property.  This is appropriate near degeneracies.

\paragraph{Metric endpoint.}
Return a rational ball, interval, finite net, or pairwise diameter bound.  This exposes exactly the quantitative stability that is proved.

\paragraph{Residual endpoint.}
Return a theorem conditional on a producer formula.  The finite checker and the deep producer remain separate modules.

A useful framework must preserve the distinction among these modes all the way to the final object-language type.

\subsection{The embedding contract}

For each fixed theorem and query, a successful SCOPE embedding supplies:

\begin{enumerate}
\item explicit codes and bounds;
\item fixed primitive-recursive algorithms;
\item characteristic bits for all bounded conditions;
\item quantifier-free loop invariants;
\item object derivations for local steps;
\item a justified certificate escape mode;
\item exact residual formulas for unavailable producers;
\item a final composition producing \(\Prf\,\PRAAx\).
\end{enumerate}

This contract is deliberately stricter than ``the algorithm terminates'' or ``Lean proves the semantic proposition.''  It is what makes the result portable, auditable, and reusable by other iPRA developments.

\section{Scoped contracts, residuals, and certificate escape}
\label{sec:scope-contracts}

\subsection{Structural lemmas}

We first isolate the proof-theoretic operations.  They are simple, but making them explicit prevents later ambiguity about where a witness lives.

\begin{lemma}[Boundary introduction]
If \(B,\Gamma\vdash_{\PRAAx} A\), then \(\Gamma\vdash_{\PRAAx} B\to A\).
\end{lemma}
\begin{proof}
Apply implication introduction.  This is exactly the \code{Prf.impI} constructor.
\end{proof}

\begin{lemma}[Residual discharge]
If \(\Gamma\vdash_{\PRAAx} B\) and \(B,\Gamma\vdash_{\PRAAx} A\), then \(\Gamma\vdash_{\PRAAx}A\).
\end{lemma}
\begin{proof}
Introduce the conditional by the preceding lemma and apply implication elimination to the proof of \(B\).
\end{proof}

\begin{lemma}[Simultaneous residual cut]
Let \(R=[B_1,\ldots,B_r]\).  If \(R,\Gamma\vdash_{\PRAAx}A\), and for every \(B_i\in R\) there is a derivation \(\Gamma\vdash_{\PRAAx}B_i\), then \(\Gamma\vdash_{\PRAAx}A\).
\end{lemma}
\begin{proof}
Use the checked theorem \code{Prf.trans}.  For a hypothesis in \(R\), supply its assumed derivation.  For a hypothesis already in \(\Gamma\), use the hypothesis rule.  The list order is irrelevant because membership, not position, is the premise of \code{Prf.trans}.
\end{proof}

\begin{lemma}[Specialization]
If \(\Gamma(\vec x)\vdash_{\PRAAx}A(\vec x)\) and \(\vec t\) is any tuple of object terms, then
\[
  \Gamma(\vec t)\vdash_{\PRAAx}A(\vec t).
\]
\end{lemma}
\begin{proof}
Apply \code{Prf.subst}.  Capture avoidance beneath quantifiers is already established by \code{Syn.up} and the formula substitution lemmas.
\end{proof}

\subsection{Private certificate sealing}

Let \(V(c,\vec x)\) be a formula with a fresh certificate variable \(c\), and let \(A(\vec x)\) not contain \(c\).

\begin{theorem}[Private sealing]
If
\[
  \Gamma\vdash_{\PRAAx}\exists c\,V(c)
  \qquad\text{and}\qquad
  V(c),\Gamma^{+}\vdash_{\PRAAx}A^{+},
\]
then \(\Gamma\vdash_{\PRAAx}A\).  Here \({}^{+}\) denotes weakening by the fresh certificate variable.
\end{theorem}
\begin{proof}
This is the existential-elimination constructor \code{Prf.exE}.  Its scope condition is syntactic: the conclusion \(A\) lives outside the certificate scope, so the witness cannot occur in it.  No choice function is created.
\end{proof}

This theorem is the default SCOPE composition rule.  A producer can use any valid certificate internally, and a consumer can inspect it, as long as the public conclusion is independent of the witness.

\subsection{Existential output escape}

A certificate-dependent object may still escape under an existential quantifier.

\begin{theorem}[Existential output]
Suppose \(V(c)\) is produced and, under \(V(c)\), one can derive \(P(t(c))\).  Then one can derive \(\exists z\,P(z)\) without naming a global function \(c\mapsto t(c)\).
\end{theorem}
\begin{proof}
Inside the certificate scope, apply existential introduction with witness \(t(c)\), obtaining \(\exists z\,P(z)\).  This formula no longer contains the certificate variable.  Then eliminate \(c\) by private sealing.  The two existential quantifiers have different roles: the inner derivation chooses a result from a local certificate, while the outer theorem states only existence.
\end{proof}

This rule is often preferable to least search.  Canonical naming should not be introduced merely because a proof assistant can order codes.

\subsection{Stable observation escape}

Suppose the public consumer needs only \(O(c)\), not the full certificate.

\begin{definition}[Stable observer]
An observer is stable for \(V\) if iPRA proves that any two valid certificates have observations related by the public equality notion.  The relation may be literal equality, an explicit quotient congruence, or a metric bound.
\end{definition}

For literal equality, one may either prove a certificate-free term \(o(q)\) with \(V(c)\to O(c)=o(q)\), or select one bounded valid certificate internally and prove that all others have the same observation.  The public theorem exposes only \(o\), not the selected certificate.  Quotient and metric variants are treated in \cref{sec:observers}.

\subsection{Canonical escape}

A public theorem sometimes genuinely needs a concrete code: a least counterexample, a normalized finite presentation, or a reproducible table.  Then a bounded order supplies a constructive canonical choice.

\begin{definition}[Canonical bounded certificate]
Assume \(c<C(q)\) and \(\Valid(q,c)\) is decidable.  Define
\[
  \mu(q)=\min\{c<C(q):\Valid(q,c)\},
\]
with \(C(q)\) itself as the failure sentinel.
\end{definition}

The checked bounded-search library proves the soundness and minimality properties needed for this definition.  The producer \(\exists c<C\,\Valid(c)\) implies \(\mu<C\) and \(\Valid(\mu)\).  Canonical escape is therefore not an application of unbounded choice; it is finite least search with an explicit guard.

\subsection{Residual escape}

If a producer cannot be constructed, SCOPE does not insert a fictitious certificate.  The producer formula itself becomes a boundary.  For example, if a finite enumeration needs a height \(H\), the theorem may have the form
\[
  \mathsf{AllSolutionsBelow}(H),\Gamma\ \vdash_{\PRAAx}\
  \mathsf{CompleteList}(\mathsf{enumerate}(H)).
\]
This is already useful: later mathematics need prove only the height formula, while the enumeration and completeness checker are generic.

\subsection{Escape effects and composition}

Assign each certificate-producing node one of the effects
\[
  \mathsf{Private},\quad \mathsf{Existential},\quad \mathsf{Stable},
  \quad \mathsf{Canonical},\quad \mathsf{Residual}.
\]
The effect is determined by the public dependency of downstream nodes, not by how the certificate was discovered.

Composition follows these rules:

\begin{enumerate}
\item a private certificate is eliminated as soon as its last consumer has been compiled;
\item an existential output may depend on the certificate only in the witness term under the public existential;
\item a stable output must pass through a proved observer interface;
\item a canonical output must have a finite bound, decidable validity bit, and least-search proof;
\item a residual output is added to the residual context and never treated as produced data.
\end{enumerate}

The rules are local and syntax-directed.  They amount to a small effect system for finite mathematical evidence.

\subsection{Scope normal form}

We can now state the central composition theorem.

\begin{theorem}[SCOPE normal form]
\label{thm:scope-normal-form}
Fix a finite interaction graph \(G\) and a finite query \(q\).  Assume:

\begin{enumerate}
\item every deterministic node is a fixed primitive-recursive transducer with an object-theory derivation of its local specification;
\item every bounded predicate has a reflected characteristic bit;
\item every loop has a quantifier-free invariant with object-theory base and step derivations;
\item every cyclic demand component is monotone on an explicitly capped finite product;
\item every relational node has finite validity and compatibility bounds;
\item every certificate edge is assigned one of the five escape modes, with the corresponding local proof;
\item every unavailable producer leaf is labeled by an object formula \(B_i(q)\).
\end{enumerate}

Then the graph can be compiled, at the meta-level, to a finite list of reachable residuals \(R(q)=[B_1(q),\ldots,B_r(q)]\) and an object derivation
\[
  R(q),\Gamma(q)\ \vdash_{\PRAAx}\ \Goal_G(q).
\]
No internal certificate occurs free in the public goal.  If every residual is derivable from a context \(\Delta\), then \(\Delta\vdash_{\PRAAx}\Goal_G(q)\).
\end{theorem}

\begin{proof}
First close every cyclic demand component by the finite theorem of \cref{sec:demand-closure}.  This produces concrete resource terms and converts the remaining data graph into an acyclic, query-specialized graph.

Topologically order the nodes.  Compile an exact transducer by specializing its local derivation and composing it with the derivations of its inputs.  Compile a Boolean or bounded-search node by the reflection theorems of \cref{sec:bounded-reflection}.  Compile a loop by fixed-trace reflection from \cref{sec:trace}.  Compile a relational path node by backward viability from \cref{sec:viability}.

At a certificate boundary, apply the declared escape rule.  For a private certificate, use existential elimination immediately after its last consumer.  For an existential output, introduce the public existential inside the certificate scope and then eliminate the certificate.  For a stable output, apply its observer theorem and eliminate the certificate.  For a canonical output, instantiate the consumer with the least accepted term.  For a residual, add the producer formula to the local context.

Because the graph is finite, this process terminates.  At each step the compiler holds an actual \(\Prf\,\PRAAx\) value.  Context mismatches are repaired by checked weakening and substitution.  Local derivations are joined by checked simultaneous cut.  All internal certificate variables are discharged at their lexical scopes, so none can occur in the outer goal.

The residual list is the dependency-normal set of producer leaves reachable from the goal.  Leaves not used by the compiled subgraph do not appear.  Finally, if a context \(\Delta\) derives every residual and every original assumption required by \(\Gamma\), apply simultaneous cut once more to obtain the closed or differently contextualized theorem.
\end{proof}

\begin{remark}
The theorem is a compilation schema for each fixed graph.  It does not quantify inside iPRA over arbitrary graph codes or arbitrary primitive-recursive programs.  This is why it is compatible with the absence of a universal primitive-recursive evaluator.
\end{remark}

\subsection{Why scope normalization resolves hidden choice}

Hidden choice occurs when a proof establishes \(\forall p\,\exists c\,V(p,c)\) and then silently treats the independently chosen \(c\)'s as one coherent global name.  SCOPE blocks this move.  At each finite query the certificate is private unless a stronger interface is proved.  A coherent name requires a canonical rule compatible across precisions, or a shrinking-diameter and motion modulus as in \cref{sec:quantitative-presentations}.  Without that data, the correct endpoint remains existential, relational, metric, or residual.

\section{Fixed bounded-formula reflection}
\label{sec:bounded-reflection}

Bounded formulas are the main bridge between ordinary finite reasoning and quantifier-free induction.  The compiler described here is a schema for each fixed formula.  It is not an object-level interpreter for arbitrary syntax codes.

\subsection{The bounded syntax}

Fix free variables \(\vec x\).  Let \(\mathcal B\) be generated by:
\[
\begin{aligned}
A ::= {}& s(\vec x)=t(\vec x) \mid \bot \mid A\wedge A \mid A\vee A
\mid A\to A \\
&\mid \exists y<B(\vec x)\,A(y,\vec x)
\mid \forall y<B(\vec x)\,A(y,\vec x),
\end{aligned}
\]
where every term denotes a fixed primitive-recursive function.  The bound is independent of the newly bound variable.  In the actual first-order syntax, the two bounded quantifiers abbreviate
\[
\exists y\,(y<B\wedge A),
\qquad
\forall y\,(y<B\to A).
\]

Implication is included because application theorems naturally use it.  Its Boolean compilation is justified constructively because every formula in \(\mathcal B\) will be proved decidable.

\subsection{Boolean terms}

A term \(b\) is a \emph{proved bit} when iPRA derives \(b=0\vee b=1\).  For proved bits define
\[
\begin{aligned}
\neg_b b &= 1\mathbin{\dot-}b,\\
b\wedge_b c &= bc,\\
b\vee_b c &= 1\mathbin{\dot-}\bigl((1\mathbin{\dot-}b)(1\mathbin{\dot-}c)\bigr),\\
b\Rightarrow_b c &= (\neg_b b)\vee_b c.
\end{aligned}
\]
Here \(\dot-\) is truncated subtraction.  All operations are primitive recursive.

\begin{lemma}[Boolean truth tables]
For proved bits \(b,c\), iPRA derives that the displayed operations are bits and satisfy the usual truth tables.  In particular,
\[
\begin{aligned}
 b\wedge_b c=1 &\leftrightarrow b=1\wedge c=1,\\
 b\vee_b c=1 &\leftrightarrow b=1\vee c=1,\\
 \neg_b b=1 &\leftrightarrow b=0,\\
 b\Rightarrow_b c=1 &\leftrightarrow (b=0\vee c=1).
\end{aligned}
\]
\end{lemma}
\begin{proof}
Use the two supplied disjunctions \(b=0\vee b=1\) and \(c=0\vee c=1\).  There are four finite cases.  In each case the result follows from the checked defining equations and semiring/monus laws.  This is intuitionistic case elimination on proved finite alternatives, not unrestricted excluded middle.
\end{proof}

Atomic equality is compiled by the checked term \(\mathsf{eqb}(s,t)\).  The current library proves \(s=t\to\mathsf{eqb}(s,t)=1\) and \(s\ne t\to\mathsf{eqb}(s,t)=0\).  Combined with the checked decidability of term equality, this gives the atomic reflection package.

\subsection{Arbitrary-arity finite folds}

Let \(b(i,\vec x)\) be a fixed bit term.  Define
\[
\begin{aligned}
\mathsf{All}_b(0,\vec x)&=1,&
\mathsf{All}_b(n+1,\vec x)&=\mathsf{All}_b(n,\vec x)b(n,\vec x),\\
\mathsf{Any}_b(0,\vec x)&=0,&
\mathsf{Any}_b(n+1,\vec x)&=\mathsf{Any}_b(n,\vec x)\vee_b b(n,\vec x).
\end{aligned}
\]

These definitions work with any fixed number of parameters.  The implementation uses \code{PTm.toPR} to abstract the open body and \code{recDef} to introduce the recursion argument and accumulator.  The parameter-renaming map sends the body's index variable to step variable \(0\), skips the accumulator at step variable \(1\), and sends every outer parameter two positions upward.  This construction is written explicitly in \code{lean/SCOPE/BoundedFormulaBlueprint.lean}.

\begin{lemma}[Finite existential reflection]
\label{lem:any-reflection}
For a proved bit body,
\[
  \mathsf{Any}_b(n,\vec x)=1
  \quad\leftrightarrow\quad
  \exists i<n\; b(i,\vec x)=1.
\]
Moreover a least such \(i\) is primitive recursive, with \(n\) as the failure sentinel.
\end{lemma}
\begin{proof}
Search for a zero of the complement \(1\mathbin{\dot-}b(i,\vec x)\).  The checked search term \(\mu(n,\vec x)\) satisfies \(\mu\le n\); if \(\mu<n\), then the complement at \(\mu\) is zero, hence \(b(\mu)=1\); and every earlier index has nonzero complement, hence bit value zero for the original body is impossible there.  If a witness \(i<n\) with \(b(i)=1\) is supplied, \code{search_found} yields \(\mu<n\), validity at \(\mu\), and \(\mu\le i\).

It remains to connect the recursive disjunction fold to the search success bit.  Prove by quantifier-free induction on \(n\) the equation
\[
  \mathsf{Any}_b(n)=1
  \quad\leftrightarrow\quad
  \mu(n)<n.
\]
The induction matrix can be expressed as an equality of characteristic terms.  At the successor step, either a witness was already found, in which case both sides remain true, or the sentinel was \(n\); then both sides become true exactly when \(b(n)=1\).  All branching is on reflected term comparisons and body bits.  Combining the equation with the checked search specification gives the desired existential equivalence and leastness.
\end{proof}

\begin{lemma}[Finite universal reflection]
\label{lem:all-reflection}
For a proved bit body,
\[
  \mathsf{All}_b(n,\vec x)=1
  \quad\leftrightarrow\quad
  \forall i<n\;b(i,\vec x)=1.
\]
\end{lemma}
\begin{proof}
Search for a counterexample, that is, an index with \(b(i)=0\).  Equivalently search for a zero of \(b\) itself.  Let \(\nu(n)\) be the least zero or sentinel \(n\).  A quantifier-free induction on the recursion equations proves
\[
  \mathsf{All}_b(n)=1 \quad\leftrightarrow\quad \nu(n)=n.
\]

For the forward implication, suppose \(i<n\).  If \(b(i)=0\), checked search foundness gives \(\nu<n\), contradicting \(\nu=n\).  Since \(b(i)\) is a bit, it must equal one.

For the backward implication, assume every \(i<n\) has value one.  If \(\nu<n\), the search specification gives \(b(\nu)=0\), while the universal premise instantiated at \(\nu\) gives \(b(\nu)=1\), a contradiction.  Therefore \(\nu=n\), and the fold equation yields \(\mathsf{All}_b(n)=1\).

No induction predicate contains the bounded universal formula.  The only induction is on equations defining the fold and search state.
\end{proof}

\subsection{The reflection theorem}

\begin{theorem}[Fixed bounded reflection]
\label{thm:bounded-reflection}
For every fixed \(A(\vec x)\in\mathcal B\), there is a primitive-recursive term \(\bit_A(\vec x)\) and iPRA derivations, uniform in the free parameters, of
\[
\begin{aligned}
&\bit_A=0\vee\bit_A=1,\\
&A\to\bit_A=1,\\
&\bit_A=1\to A.
\end{aligned}
\]
Consequently iPRA proves \(A\leftrightarrow\bit_A=1\) and \(A\vee\neg A\).
\end{theorem}

\begin{proof}
Proceed by structural recursion on the fixed syntax tree.

For an equality \(s=t\), take \(\bit_A=\mathsf{eqb}(s,t)\).  Equality decidability and the two checked \code{eqb} lemmas provide Booleanity and both implications.  For \(\bot\), take the constant zero bit.  The forward implication follows by bottom elimination and the backward implication is \(0=1\to\bot\), using successor separation.

For conjunction and disjunction, apply the induction hypotheses and the Boolean truth-table lemma.  For example, from \(A\wedge B\) obtain both component bits equal to one and hence their product equal to one.  Conversely, product one implies both bits one; reflect them back and introduce the conjunction.

For implication, let \(a=\bit_A\) and \(b=\bit_B\).  Suppose \(A\to B\).  Decide \(A\) using the induction hypothesis.  If \(\neg A\), then \(a=0\) and the implication bit is one.  If \(A\), obtain \(B\), hence \(b=1\), and again the implication bit is one.  Conversely, if \(a\Rightarrow_b b=1\), the truth table gives \(a=0\vee b=1\).  Under an assumption \(A\), the first alternative contradicts the forward reflection of \(A\); the second reflects to \(B\).  Thus \(A\to B\) follows intuitionistically.

For \(\exists y<B\,A(y)\), compile the body to \(b(y)\) and take \(\bit=\mathsf{Any}_b(B)\).  A supplied existential witness gives \(i<B\) and \(A(i)\), hence \(b(i)=1\), so \cref{lem:any-reflection} gives the aggregate bit.  Conversely, aggregate bit one supplies the least \(i<B\) with \(b(i)=1\); the induction hypothesis yields \(A(i)\), and existential introduction finishes.

For \(\forall y<B\,A(y)\), take \(\mathsf{All}_b(B)\).  A universal proof gives each body bit one, hence the fold bit by \cref{lem:all-reflection}.  Conversely, fold bit one gives \(b(i)=1\) for an arbitrary \(i<B\); reflect to \(A(i)\) and universally introduce the bound variable.

Booleanity is preserved at every constructor.  Finally, decide whether \(\bit_A=1\).  In the positive case reflect to \(A\).  In the negative case, Booleanity forces \(\bit_A=0\); if \(A\) held, forward reflection would give \(\bit_A=1\), contradiction.  Therefore \(A\vee\neg A\).
\end{proof}

\subsection{Binder-safe specialization}

The theorem is stated over open formulas, so substitution of arbitrary terms is obtained directly from \code{Prf.subst}.  Beneath a quantifier, \code{Syn.up} preserves the newly bound variable and shifts the substituted outer terms.  The checked identities \code{Fm.shift_subst_up} and \code{Fm.subst_inst} prove that weakening and instantiation commute in the required form.

Thus the first implementation milestone does not need a bespoke semantic substitution proof for every application.  It needs a syntactic compiler coherence theorem only to show that regenerated code is extensionally equal to the substituted old code.  The already checked \code{toPR_eq} theorem supplies the key step for terms abstracted into primitive-recursive symbols.

\subsection{Canonical guarded bounded choice}

\begin{theorem}[Least accepted witness]
\label{thm:least-accepted}
Let \(R(\vec x,y)\in\mathcal B\), and let \(B(\vec x)\) be a primitive-recursive bound.  There are primitive-recursive terms \(w(\vec x)\) and \(h(\vec x)\in\{0,1\}\) such that iPRA proves
\[
\begin{aligned}
 h=1 &\leftrightarrow \exists y<B\;R(\vec x,y),\\
 h=1 &\to w<B\wedge R(\vec x,w),\\
 h=0 &\to w=B,\\
 h=1\wedge z<w &\to \neg R(\vec x,z).
\end{aligned}
\]
\end{theorem}
\begin{proof}
Compile \(R\) to its bit \(r(y,\vec x)\).  Apply the checked least-zero search to \(1\mathbin{\dot-}r\).  Let \(w\) be the search result and let \(h\) be the reflected comparison bit for \(w<B\).  The search specification gives the sentinel law, soundness, and minimality.  The reflection theorem translates between \(r=1\) and \(R\).  Every displayed claim follows by finite intuitionistic reasoning.
\end{proof}

\begin{corollary}[Bounded choice schema]
If iPRA proves \(\forall\vec x\,\exists y<B(\vec x)\,R(\vec x,y)\), then it proves \(\forall\vec x\,R(\vec x,w(\vec x))\) for the concrete primitive-recursive least witness term \(w\).
\end{corollary}

The bound is essential.  This is not countable choice and does not turn arbitrary total relations into primitive-recursive functions.

\subsection{Why bounded reflection is enough for many classical finite arguments}

A finite decidable argument may be discovered classically, by SAT solving, computer algebra, or exhaustive search.  To import it constructively, expose a finite certificate and a bounded verifier.  The certificate need not be canonical unless downstream syntax needs its code.  In the common case, private sealing eliminates it after the verifier derives the certificate-independent theorem.

This separation permits aggressive external discovery while retaining a small trusted object theory: only the verifier equations and their \(\Prf\,\PRAAx\) correctness proofs matter to the final theorem.

\section{Fixed-trace reflection and proof-producing finite programs}
\label{sec:trace}

\subsection{State machines specialized at the theorem}

Fix a query tuple \(q\), a state code \(s\), and a transition function
\[
  T(s,q).
\]
All are represented by fixed primitive-recursive symbols.  Define the trace by primitive recursion:
\[
  \mathsf{run}(0,q)=s_0(q),
  \qquad
  \mathsf{run}(t+1,q)=T(\mathsf{run}(t,q),q).
\]
The current library's \code{recDef} gives the symbol and object-theory zero/successor equations at arbitrary terms.

A state may pack several logical components: loop variables, accumulators, a failure bit, table cursors, certificate digits, and resource counters.  Pairing or mixed-radix coding turns the tuple into one natural number.  Fixed-dimensional tuples do not require arrays as primitive objects.

\subsection{Quantifier-free invariants}

Let \(I(s,q)\) be a quantifier-free formula or, preferably, an invariant bit equation \(\iota(s,q)=1\).  Assume object derivations
\[
\begin{aligned}
&\Gamma(q)\vdash_{\PRAAx} I(s_0(q),q),\\
&I(s,q),\Gamma(q)^{+}\vdash_{\PRAAx} I(T(s,q),q).
\end{aligned}
\]
The fresh variable in the step sequent is the induction counter/state parameter as arranged by \code{IPRA.ind}.

\begin{theorem}[Fixed-trace reflection]
\label{thm:trace-reflection}
For every fuel term \(N(q)\), iPRA derives
\[
  \Gamma(q)\vdash_{\PRAAx} I(\mathsf{run}(N(q),q),q).
\]
If an object derivation extracts a terminal formula \(F(q)\) from that invariant, then \(\Gamma(q)\vdash_{\PRAAx}F(q)\).
\end{theorem}
\begin{proof}
Form the open induction formula
\[
  \phi(t,q):=I(\mathsf{run}(t,q),q).
\]
It is quantifier-free because \(I\) is and \(\mathsf{run}\) is a term.  The base case reduces by the zero equation of \(\mathsf{run}\) to the supplied base derivation.  The successor case reduces by the successor equation to the supplied step derivation.  Apply \code{IPRA.ind}, then instantiate the induction variable with \(N(q)\) by \code{Prf.subst}.  Compose the terminal extractor by simultaneous cut or implication elimination.
\end{proof}

\subsection{Booleanizing prefix properties}

An informal loop invariant often says ``all entries processed so far satisfy \(P\).''  The literal formula \(\forall i<t\,P(i)\) is not quantifier-free.  SCOPE instead extends the state with
\[
  a(0)=1,
  \qquad
  a(t+1)=a(t)\wedge_b \bit_P(t).
\]
The invariant contains the equation \(a(t)=1\).  After induction, bounded reflection identifies it with the intended prefix property.  A failure accumulator is dual:
\[
  f(0)=0,
  \qquad
  f(t+1)=f(t)\vee_b\neg_b\bit_P(t).
\]

This pattern handles table verification, finite covers, path checks, proof traces, and bounded universal side conditions.  It also makes counterexamples extractable by least search.

\subsection{Adaptive finite interactions}

A consumer may ask its next finite question based on previous answers.  Such adaptivity is still a fixed trace if:

\begin{enumerate}
\item the interaction horizon is bounded by a primitive-recursive term;
\item every answer lies in an explicit finite code range;
\item the next-query function is primitive recursive;
\item validity of each answer is reflected by a bit.
\end{enumerate}

The state stores the transcript.  The transition decodes the current answer, verifies it, updates the transcript and next query, and accumulates failure.  A private master certificate can encode the entire answer table in advance; the trace merely reads the demanded row.  This ``offline certificate, online verifier'' view is often easier to formalize than a genuinely interactive object.

\subsection{Phase composition}

Large algorithms are best split into phases with separate invariants.  Suppose phase \(j\) proves a postcondition \(P_j\), and phase \(j+1\) has a base derivation from \(P_j\).  Compile each phase by \cref{thm:trace-reflection} and compose the resulting derivations with \code{Prf.trans}.  There is no need for one enormous invariant carrying all earlier implementation details.

Phases also localize certificate scopes.  A certificate used only to establish \(P_j\) can be existentially eliminated before phase \(j+1\), preventing irrelevant data from polluting the public theorem and subsequent state.

\subsection{Termination and descent}

A loop whose natural description is ``repeat until done'' must expose primitive-recursive fuel.  Typical sources are:

\begin{itemize}
\item a finite search bound;
\item a decreasing natural rank;
\item a lexicographic tuple encoded by a rank;
\item a finite state-space cardinality together with a no-repeat certificate;
\item the capped demand budget of \cref{sec:demand-closure}.
\end{itemize}

The repository already contains generic descent developments.  SCOPE uses them when the step relation is theorem-specific, but retains the same discipline: termination evidence is a finite term and its correctness is an object derivation.

\subsection{A restricted proof-producing DSL}

The fixed-trace theorem suggests a small implementation language with nodes for:

\begin{itemize}
\item exact term computation;
\item reflected conditionals;
\item bounded folds and least search;
\item clocked recursion with an invariant;
\item certificate read/verify operations;
\item local existential scopes;
\item residual producer calls.
\end{itemize}

A compiler for this DSL need not be trusted for logical correctness.  It emits \(\PR\) symbols, formulas, and a \(\Prf\,\PRAAx\) term assembled from checked lemmas.  Lean checks the output.  The DSL is deliberately weaker than arbitrary Lean, which is what makes proof production feasible.

\section{Capped cyclic demand closure}
\label{sec:demand-closure}

Many quantitative proofs have backward dependencies.  An output precision asks one input for more precision; that input's construction asks another component for more; and a cycle may return to the first component.  Running such a cycle dynamically obscures the induction bound.  SCOPE closes all demands first in a finite lattice.

\subsection{The finite demand lattice}

Fix \(k\) coordinates and caps \(C_0,\ldots,C_{k-1}\).  Let
\[
  D=\prod_{i<k}\{0,1,\ldots,C_i\}
\]
with coordinatewise order.  Let \(F:D\to D\) be monotone.  Define the inflationary transformer
\[
  H(d)=d\vee F(d),
\]
where the join is coordinatewise maximum.  Starting from \(d_0=0\), put \(d_{t+1}=H(d_t)\).

\begin{lemma}[Monotone bounded iteration]
The sequence \((d_t)\) is nondecreasing and remains in \(D\).
\end{lemma}
\begin{proof}
By definition, \(d_t\le d_t\vee F(d_t)=d_{t+1}\).  Both operands of the join are in \(D\), so each coordinate remains below its cap.  Induction on \(t\) proves the claim.
\end{proof}

Define
\[
  \rank(d)=\sum_{i<k}d_i,
  \qquad
  K=\sum_{i<k}C_i.
\]
Then \(0\le\rank(d)\le K\).

\begin{lemma}[Strict rank increase]
If \(H(d)\ne d\), then \(\rank(H(d))\ge\rank(d)+1\).
\end{lemma}
\begin{proof}
Since \(d\le H(d)\) and the vectors are unequal, at least one coordinate strictly increases.  Natural coordinates are integral, so that coordinate increases by at least one.  Every other coordinate is nondecreasing.  Summing gives the inequality.
\end{proof}

\begin{theorem}[Finite cyclic closure]
\label{thm:demand-closure}
The iterate \(d_K\) is a fixed point of \(H\).  It is the least prefixed point of \(F\), equivalently the least fixed point of \(H\):
\[
 H(d_K)=d_K,
 \qquad
 F(e)\le e\Longrightarrow d_K\le e.
\]
\end{theorem}
\begin{proof}
Suppose \(d_K\) were not fixed.  Once an iterate is fixed, every later iterate is the same, so none of \(d_0,\ldots,d_K\) could have been fixed.  By strict rank increase at each of the \(K+1\) transitions from \(d_0\) through \(d_{K+1}\), and \(\rank(d_0)=0\), we would have \(\rank(d_{K+1})\ge K+1\).  This contradicts the bound \(\rank(d_{K+1})\le K\).  Hence \(d_K\) is fixed.

For leastness, let \(e\) satisfy \(F(e)\le e\).  Then \(H(e)=e\).  We prove \(d_t\le e\) by induction.  It holds at zero.  If \(d_t\le e\), monotonicity of \(F\) gives \(F(d_t)\le F(e)\le e\), and also \(d_t\le e\); therefore their join \(d_{t+1}\le e\).  In particular \(d_K\le e\).
\end{proof}

\subsection{Query-dependent caps}

The caps may depend on the finite query \(q\).  Then
\[
  K(q)=\sum_{i<k}C_i(q)
\]
is primitive recursive, and the same proof is uniform in \(q\).  The dimension \(k\) is fixed by the theorem graph; it is not an object-level unbounded parameter unless a separate coding theorem is supplied.

The cap is part of the contract.  A demand transformer without a cap does not automatically have a primitive-recursive fixed-point modulus.  In applications the cap comes from the terminal precision, a maximum finite search height, an explicit degree bound, or a finite resource policy.

\subsection{Object-theory arithmetization}

For a fixed \(k\), encode vectors by mixed radix with bases \(C_i+1\).  Digit extraction is implemented by quotient and remainder.  The checked division library supplies the defining equations and uniqueness properties.  Alternatively, for small \(k\), nested pairing may be simpler.

The object proof needs the following reflected terms:

\begin{itemize}
\item \(\mathsf{digit}_i(d,q)\);
\item the code of \(F(d,q)\) and \(H(d,q)\);
\item \(\rank(d,q)\) and \(K(q)\);
\item a fixed-point bit \(\mathsf{eqb}(H(d),d)\);
\item coordinate-bound and monotonicity bits.
\end{itemize}

The central induction invariant is quantifier-free.  It states that the trace code is the prescribed iterate, each digit is bounded, rank is bounded, and a ``seen fixed point'' bit is persistent.  Strict rank increase is a local arithmetic implication derived from reflected coordinate comparison.  The contradiction after \(K+1\) nonfixed steps is an inequality of natural terms.

Thus the human proof above translates directly to \code{IPRA.ind}; no induction over a completed lattice or an object-language universal quantifier is required.

\subsection{Acyclic and cyclic demand compilation}

For an acyclic graph, demands are propagated once in reverse topological order.  Each node has a primitive-recursive transformer from output requests to input requests.  Taking coordinatewise maxima combines multiple consumers.

For a cyclic strongly connected component, collect its local transformers into \(F\), apply \cref{thm:demand-closure}, and replace every internal demand by the resulting fixed vector.  Collapse the component to one acyclic supernode.  Repeating this for all strongly connected components yields a global finite schedule.

\subsection{Correctness of reverse demand}

A demand transformer is meaningful only with a local theorem of the form:
\begin{quote}
If every input is valid at the demanded resources computed from output request \(p\), then the node's output is valid at \(p\).
\end{quote}

At a fixed point, every cyclic edge receives at least the resource requested by its consumer.  Compose the local theorems around the component, then run the forward trace once at the closed resources.  Demand closure therefore separates termination/scheduling from numerical correctness.

\subsection{Why monotonicity is essential}

If a transformer is not monotone, iteration from zero need not compute a least solution, and increasing one resource could perversely reduce another requirement.  Such a contract is unsuitable for modular quantitative reasoning.  The remedy is usually to replace exact demands by monotone upper envelopes.  This may over-approximate resources, but the proof remains finite and compositional.

The executable reference implementation checks monotonicity exhaustively for small examples and records every iterate.  In Lean, monotonicity is an object derivation or a host theorem used to construct the object derivation for the fixed transformer.

\section{Master certificates, canonical naming, and finite strategies}

\subsection{Packing finite evidence into one code}

A fixed proof graph may need many finite pieces of evidence: table rows, path nodes, local witnesses, branch flags, and trace annotations.  Suppose component \(j\) has range \(0\le c_j<R_j(q)\).  Define the mixed-radix code
\[
  c=c_0+R_0c_1+R_0R_1c_2+\cdots.
\]
The total capacity is \(C(q)=\prod_jR_j(q)\).  Quotient and remainder recover every digit.

\begin{lemma}[Mixed-radix correctness]
For fixed finitely many radices \(R_j\ge1\), encoding maps the product of digit ranges injectively into \([0,C)\), and decoding is its inverse.  All functions and bounds are primitive recursive, and their equations are provable in iPRA.
\end{lemma}
\begin{proof}
The arithmetic proof is the standard repeated quotient--remainder argument.  For one radix, \(c=R q+r\) with \(r<R\) has unique \(q,r\).  Iterate this decomposition through the fixed list of radices.  The checked division theorem supplies existence and uniqueness at every step; finite induction over the fixed number of components can be expanded at the meta-level or represented by a bounded trace.
\end{proof}

A \emph{master certificate} is such a code together with a verifier bit \(V(q,c)\).  The verifier decodes only the finite fields required by the graph and checks local relations.  There is no semantic need for the certificate format to mirror the source mathematics exactly; its purpose is to support a small object verifier.

\subsection{One verifier, many local contracts}

Let \(V_j(q,c)\) be the bit for the \(j\)-th local contract.  Define
\[
  V(q,c)=\bigwedge_j V_j(q,c)
\]
by a fixed finite conjunction or a Boolean fold.  Bounded reflection proves that \(V=1\) is equivalent to every checked local condition.  The extractor can then use the decoded fields under one certificate scope.

Packing evidence has three benefits:

\begin{enumerate}
\item one existential producer formula replaces a long nest of dependent witnesses;
\item one least bounded search can canonically name the entire evidence package if necessary;
\item the final theorem can keep the certificate private and expose only stable observations.
\end{enumerate}

It does not force canonical naming.  The default remains private existential elimination.

\subsection{Least accepted master certificate}

\begin{theorem}[Canonical master certificate]
Let \(V(q,c)\) be a reflected validity bit and \(C(q)\) an explicit bound.  There is a primitive-recursive term \(\mu(q)\) such that iPRA proves:
\[
\begin{aligned}
&\mu\le C,\\
&\mu<C\to V(q,\mu)=1,\\
&j<\mu\to V(q,j)=0,\\
&\bigl(\exists c<C\;V(q,c)=1\bigr)\to \mu<C,\\
&\neg\exists c<C\;V(q,c)=1\to\mu=C.
\end{aligned}
\]
\end{theorem}
\begin{proof}
Apply the checked least-zero search to \(1\mathbin{\dot-}V(q,c)\).  The search term uses \(C\) as its sentinel.  Soundness, minimality, foundness, and the sentinel conclusion are exactly the checked search specification translated through Booleanity of \(V\).
\end{proof}

This theorem resolves reproducibility requirements without adding a global choice principle.  It should be used only when the public theorem needs the code or data derived from a canonical code.  If the goal is certificate-independent, private sealing is smaller and mathematically cleaner.

\subsection{Finite bounded games}

Nested bounded quantifiers can be evaluated as a finite game.  Consider a fixed prefix
\[
  Q_0x_0<B_0\;Q_1x_1<B_1(x_0)\;\cdots\;
  Q_{d-1}x_{d-1}<B_{d-1}(\vec x_{<d-1})\;R(\vec x),
\]
where each \(Q_i\) is \(\exists\) or \(\forall\), and \(R\) is decidable.  At existential nodes the Proponent chooses a child; at universal nodes the Opponent chooses.

Define the winning bit backward from the leaves:
\[
\begin{aligned}
W_d(\vec x)&=\bit_R(\vec x),\\
W_i(\vec x_{<i})&=
\begin{cases}
\mathsf{Any}_{x_i<B_i} W_{i+1}(\vec x_{\le i}),&Q_i=\exists,\\
\mathsf{All}_{x_i<B_i} W_{i+1}(\vec x_{\le i}),&Q_i=\forall.
\end{cases}
\end{aligned}
\]
Because the prefix depth is fixed, these are concrete primitive-recursive terms.

\begin{definition}[Finite strategy certificate]
A Proponent strategy records one child at each reachable existential history.  An Opponent counterstrategy records one child at each reachable universal history.  A table is encoded in mixed radix, and a verifier checks bounds, table lookup, and every terminal play induced by the opposing choices.
\end{definition}

\begin{theorem}[Finite strategy dichotomy]
\label{thm:finite-strategy}
For every fixed bounded game, iPRA proves a decidable dichotomy:

\begin{enumerate}
\item if \(W_0=1\), there exists a bounded Proponent strategy whose verifier is one, and every play consistent with it ends at a true leaf;
\item if \(W_0=0\), there exists a bounded Opponent counterstrategy whose verifier is one, and every play consistent with it ends at a false leaf.
\end{enumerate}
\end{theorem}
\begin{proof}
Induct backward on the fixed game depth in the metatheory while constructing object derivations.

At a leaf, the reflected matrix bit decides which trivial certificate is valid.  At an existential node with \(W_i=1\), finite existential reflection supplies the least winning child.  Store that child and recursively store the winning strategy below it.  If \(W_i=0\), every child has losing bit zero; collect the recursively constructed counterstrategies for all children in one bounded table.

At a universal node the roles reverse.  If \(W_i=1\), every child is winning, so collect a winning strategy for each child.  If \(W_i=0\), finite universal reflection supplies a least losing child; store it and its counterstrategy.

All collections are finite and bounded.  Soundness of a strategy verifier is proved by forward induction on the fixed remaining depth, reading the appropriate table row at each history.  Completeness follows from the backward construction.  The two alternatives are decidable because \(W_0\) is a bit.
\end{proof}

For the common shape \(\forall x<B\,\exists y<C(x)\,R(x,y)\), a winning strategy is simply a finite table \(\sigma(x)\).  The reference implementation constructs the least valid \(y\) for each \(x\) and verifies the table.

\subsection{Strategies versus formulas}

The characteristic bit of \cref{thm:bounded-reflection} is enough to prove a bounded formula.  A strategy certificate is useful when downstream computation needs the choices or when a compact externally generated certificate should be checked.  SCOPE therefore treats strategy extraction as an optional strengthening, not as the definition of bounded truth.

\subsection{Producer boundaries}

A master verifier can be fully generic while its producer is difficult.  For example:

\begin{itemize}
\item a SAT solver may produce a finite combinatorial certificate;
\item a computer algebra system may produce a factorization or Nullstellensatz witness;
\item an analytic argument may produce a rational contour bound;
\item a Diophantine theorem may produce a height cutoff.
\end{itemize}

The SCOPE endpoint may retain \(\exists c<C\;V(q,c)=1\) as a residual.  This is often the correct modular theorem: it proves that any valid finite producer suffices, without claiming a generic method for finding one.

\subsection{Proof-size considerations}

Expanding a large master verifier into propositional syntax can create huge proof terms.  The preferred implementation shares fixed reflected lemmas:

\begin{itemize}
\item one theorem for each Boolean operation;
\item one theorem for digit extraction and table lookup;
\item one trace theorem parameterized by local step derivations;
\item one least-search adapter;
\item one certificate-sealing theorem.
\end{itemize}

A future compact certificate calculus is acceptable only if its checker soundness produces an ordinary \(\Prf\,\PRAAx\) derivation.  Reducing serialization size must not enlarge the trusted logical kernel.

\section{Relational witnesses and backward viability}
\label{sec:viability}

\subsection{Why local least choice is insufficient}

Suppose a proof proceeds through finite levels of approximants.  At level \(i\) there are finitely many locally valid candidates, and compatibility constrains consecutive levels.  Choosing the least locally valid candidate can fail even when a global path exists: the least candidate may have no continuation.

This is a constructive issue, not merely an efficiency issue.  A local selector does not prove the existence of a compatible future.  SCOPE therefore computes viability backward before selecting forward.

\subsection{Finite relational system}

Fix a horizon \(h\).  At level \(i\le h\), candidates are codes \(a<M_i(q)\).  Let \(L_i(q,a)\) be the local validity bit, and for \(i<h\) let \(E_i(q,a,b)\) be the compatibility bit from level \(i\) to \(i+1\).

Define viability bits backward:
\[
\begin{aligned}
V_h(a)&=L_h(a),\\
V_i(a)&=L_i(a)\wedge_b
  \mathsf{Any}_{b<M_{i+1}}\bigl(V_{i+1}(b)\wedge_b E_i(a,b)\bigr).
\end{aligned}
\]
All parameters and bounds are primitive recursive.  Since \(h\) is a finite query or a primitive-recursive horizon, the table can be computed by a bounded trace.

\begin{theorem}[Finite path equivalence]
\label{thm:viable-path}
There exists a locally valid compatible path \(a_0,\ldots,a_h\) exactly when
\[
  \exists a_0<M_0\;V_0(a_0)=1.
\]
More generally, \(V_i(a)=1\) exactly when \(a\) extends to a valid compatible suffix through level \(h\).
\end{theorem}
\begin{proof}
Use backward induction on \(h-i\).

At the final level, \(V_h(a)=L_h(a)\), so the claim is immediate.  Assume the claim for \(i+1\).  If \(V_i(a)=1\), Boolean reflection gives \(L_i(a)=1\) and a candidate \(b<M_{i+1}\) with \(V_{i+1}(b)=1\) and \(E_i(a,b)=1\).  By the induction hypothesis, \(b\) has a valid suffix; prefixing \(a\) gives a suffix from level \(i\).

Conversely, suppose \(a\) begins a valid suffix.  Let \(b\) be its next element.  Local validity and compatibility hold, and the remaining tail makes \(V_{i+1}(b)=1\) by the induction hypothesis.  Finite existential reflection therefore makes the second conjunct of \(V_i(a)\) equal one, so \(V_i(a)=1\).
\end{proof}

\subsection{Least viable path}

\begin{theorem}[Canonical viable path]
If a compatible path exists, there is a primitive-recursive lexicographically least viable path.  It is obtained by choosing the least \(a_0\) with \(V_0(a_0)=1\), then at each level choosing the least compatible \(a_{i+1}\) with \(V_{i+1}(a_{i+1})=1\).
\end{theorem}
\begin{proof}
The initial choice exists by \cref{thm:viable-path} and is made by bounded least search.  If \(a_i\) is viable, the defining existential in \(V_i(a_i)\) supplies at least one compatible viable successor; bounded least search names the least one.  Induction on the finite horizon constructs the path.  Soundness follows from the stored validity and compatibility bits.  Lexicographic minimality follows at the first position where another path differs: the algorithm chose the least viable extension at that prefix.
\end{proof}

The word \emph{viable} is essential.  The least merely local successor need not be extendable.

\subsection{Object-theory compilation}

There are two useful implementations.

\paragraph{Explicit viability table.}
Encode all bits \(V_i(a)\) in one master certificate and verify the backward recurrence.  This is attractive when an external solver supplies the table.

\paragraph{Computed viability trace.}
Run a primitive-recursive backward dynamic program whose state contains the current level and next-level bit table.  A Boolean fold computes each row.  The invariant states that the stored suffix table equals the recurrence and that every entry is a bit.  The matrix is quantifier-free because row correctness is represented by a fold equation.

After the table is available, a forward trace performs least viable selection.  The terminal extractor proves path validity.  Every step is an instance of bounded reflection, least search, or fixed-trace reflection.

\subsection{Relational outputs without path selection}

Not every theorem needs a path.  A consumer may depend only on a stable property shared by every viable endpoint.  Then the viability table can remain private and SCOPE exports the stable observation.  This is often preferable near branching singularities.

Alternatively, the public theorem may state the existence of a compatible path.  Inside the private certificate scope, introduce the existentially quantified path code and then eliminate the certificate.  Again no global selector is required.

\subsection{Changing horizons}

For each finite horizon \(h\), the least viable path may begin differently.  It is generally invalid to identify these finite choices as prefixes of one infinite branch.  Two quantitative conditions repair the situation:

\begin{enumerate}
\item a diameter bound saying all viable observations at requested precision \(p\) lie within \(\delta(p)\), with \(\delta(p)\to0\) at a primitive-recursive rate;
\item a motion budget bounding the distance between selected observations when the horizon or precision changes, with a primitive-recursive summable tail.
\end{enumerate}

Under these conditions, the selected finite observations form a regular Cauchy name even if the underlying certificate path is not literally coherent.  The construction is developed in \cref{sec:quantitative-presentations}.  Without such data, the correct result remains a family of finite-horizon relational theorems.

\subsection{Worked dead-end pattern}

The reference example has levels
\[
  A_0=\{0,1\},\quad A_1=\{0,1\},\quad A_2=\{0\}.
\]
Edges are \(0\to0\) and \(1\to1\) from level zero, but only \(1\to0\) from level one.  Candidate \(0\) is locally valid and least at level zero, yet it is a dead end.  Backward viability marks only \(1\) viable at levels zero and one, producing the path \((1,1,0)\).

This tiny example is an acceptance test for the Lean implementation: any purported generic selector that returns \(0\) before computing viability is wrong.

\section{Observational projection and finite descent}
\label{sec:observers}

Certificates may be noncanonical while the mathematical observation of interest is canonical or quantitatively stable.  Observer interfaces permit the observation to escape without exposing the certificate.

\subsection{Exact stable observations}

\begin{theorem}[Exact observer elimination]
Let \(V(c,q)\) be a certificate relation, \(O(c,q)\) an observation term, and \(o(q)\) a certificate-free term.  If iPRA proves
\[
  \forall c\,[V(c,q)\to O(c,q)=o(q)]
\]
and a local consumer derives \(P(O(c,q),q)\) from \(V(c,q)\), then
\[
  \exists c\,V(c,q)\quad\Longrightarrow\quad P(o(q),q).
\]
\end{theorem}
\begin{proof}
Open the certificate existential.  The stability equation transports the local proof from \(P(O(c,q),q)\) to \(P(o(q),q)\) by equality substitution.  The result no longer mentions \(c\), so existential elimination closes the scope.
\end{proof}

If only pairwise stability is known,
\[
  V(c)\wedge V(c')\to O(c)=O(c'),
\]
use the least bounded valid certificate \(\mu\) internally and define \(o=O(\mu)\).  This names a canonical observation but still does not publish the certificate unless needed.

\subsection{Existential observations need no stability}

If the public conclusion is \(\exists z\,P(z)\), then stability is unnecessary.  Under a private certificate, prove \(P(O(c))\), introduce \(z=O(c)\), and eliminate \(c\).  This simple rule resolves many apparent choice problems: the mathematical theorem asked only for existence, not for a uniform selector.

\subsection{Quotient observations}

A quotient should not be introduced as an opaque classical object if the finite theorem needs only an explicit congruence.

Let \(\sim\) be a reflected equivalence relation on observation codes.  Assume every valid certificate yields an observation, and any two valid observations are congruent:
\[
  V(c)\wedge V(c')\to O(c)\sim O(c').
\]
Then one may export:

\begin{itemize}
\item an existential representative together with proofs that all valid observations are congruent to it;
\item a canonical normal form if a bounded normalizer is available;
\item any predicate proved invariant under \(\sim\).
\end{itemize}

The proof is the exact observer argument with equality transport replaced by an explicit congruence-respect theorem for the consumer.  No quotient type or quotient choice is logically necessary.

\subsection{Metric diameter observers}

Let observations be rational codes with a reflected metric \(d\).  Suppose iPRA proves
\[
  V(c)\wedge V(c')\to d(O(c),O(c'))\le\varepsilon(q).
\]

\begin{theorem}[Metric observer elimination]
From a valid certificate one can export a center \(z\) existentially such that every valid observation lies in the ball \(B(z,\varepsilon)\).  If a least bounded certificate is chosen internally, one can export the canonical center \(O(\mu)\) and the same bound.
\end{theorem}
\begin{proof}
For the existential version, open one valid certificate \(c\), choose \(z=O(c)\), and use the pairwise-diameter theorem with arbitrary \(c'\).  Then close \(c\).  For the canonical version, least search gives a fixed valid \(\mu\); instantiate pairwise diameter at \(\mu,c'\).
\end{proof}

Metric output is often the constructively correct replacement for an unstable exact choice.  It retains enough information for continuous downstream consumers.

\subsection{Gap-guarded exact labels}

Approximate data determines an exact discrete label only with a positive gap.

\begin{theorem}[Finite gap stabilization]
Let \(L< M\) range over finite labels.  Suppose exact scores \(s_L\) have a unique minimizer \(L_*\) with
\[
  s_L\ge s_{L_*}+g\qquad(L\ne L_*),
\]
where \(g>0\).  Suppose certified approximations \(\tilde s_L\) satisfy
\[
  |\tilde s_L-s_L|\le e,
  \qquad 2e<g.
\]
Then the least minimizer of the approximate scores is exactly \(L_*\).
\end{theorem}
\begin{proof}
For \(L\ne L_*\),
\[
\tilde s_L\ge s_L-e\ge s_{L_*}+g-e
> s_{L_*}+e\ge\tilde s_{L_*}.
\]
Thus \(L_*\) is the unique approximate minimizer.  All comparisons are rational inequalities and finite least search names the label.  The strict inequality \(2e<g\) is the essential producer boundary.
\end{proof}

This theorem underlies robust root labels, eigenvector normalizations, branch selection, nearest-point decisions, and sign extraction.  Without a gap, SCOPE exports an unordered or relational object instead.

\subsection{Finite observational descent on a tree}

Local data are often attached to overlapping patches.  Global gluing can be stronger than the theorem requires.  A supplied finite spanning tree is enough to compare all local observations with a root.

Let \(T\) be a rooted finite tree.  At every vertex \(v\) there is an observation \(O_v\).  For each nonroot vertex, let \(p(v)\) be its parent.

\begin{theorem}[Exact tree descent]
If every edge has a proof \(O_v=O_{p(v)}\), then every vertex observation equals the root observation.
\end{theorem}
\begin{proof}
Induct on the primitive-recursive parent-depth of \(v\).  At depth zero, \(v\) is the root.  At successor depth, compose the edge equality with the induction hypothesis for the parent.  Because the tree is finite and rooted, depth is bounded and primitive recursive from a supplied parent table.
\end{proof}

\begin{theorem}[Metric tree descent]
If every edge has
\[
  d(O_v,O_{p(v)})\le e_v,
\]
then
\[
  d(O_v,O_{\mathrm{root}})\le
  \sum_{u\text{ on the path }v\to\mathrm{root}}e_u.
\]
\end{theorem}
\begin{proof}
Use the same depth induction and the triangle inequality.  The path-error sum is computed by primitive recursion along the parent map.  The theorem exposes the accumulated finite error instead of asserting exact global gluing.
\end{proof}

Tree descent applies to local coordinate charts, finite covers, Cech-style data, patchwise polynomial identities, and distributed certificates.  If the consumer needs only a root observation or an explicit error ball, constructing an infinitary quotient or a globally selected section is unnecessary.

\subsection{Symmetric and unordered outputs}

When labels collide, stable observations are usually symmetric:

\begin{itemize}
\item coefficients rather than ordered roots;
\item an unordered multiset of root intervals;
\item elementary symmetric functions;
\item an approximate spectral projector rather than a signed eigenvector;
\item a finite relation of local branches rather than one branch name.
\end{itemize}

SCOPE treats these as first-class public interfaces.  Exact labels are a later, gap-guarded refinement rather than the default representation.

\section{Quantitative presentations of mathematical objects}
\label{sec:quantitative-presentations}

SCOPE does not prescribe one universal encoding of all mathematics.  It prescribes the interfaces that make an encoding usable in iPRA: finite queries, primitive-recursive realizers, explicit error or search bounds, and object derivations.

\subsection{Integers and rationals}

Integers may be represented as signed pairs or differences of naturals.  The current derivation library already uses a difference representation and a ring of term classes modulo provable equality.  Rational numbers can be represented by triples \((s,a,b)\) with sign \(s\), numerator \(a\), and positive denominator \(b\), together with a normalization or a cross-multiplication equality relation.

For most quantitative theorems, canonical lowest terms are optional.  Cross multiplication is a stable quotient interface.  Canonical normalization is useful only when exact code equality or least-code comparison is required.

Dyadic rationals are especially convenient for error budgets.  A dyadic \(z/2^p\) can be stored by an integer numerator and precision exponent.  Addition and comparison are primitive recursive after denominator alignment.  Absolute-value inequalities reduce to integer inequalities.

\subsection{Regular real names}

\begin{definition}[Regular name]
A regular real name is a primitive-recursive rational sequence \(a(p)\) with an explicit law such as
\[
  |a(p+1)-a(p)|\le 2^{-p}.
\]
More generally, a name may carry a primitive-recursive tail modulus \(M(p)\) satisfying
\[
  m,n\ge M(p)\Longrightarrow |a(m)-a(n)|\le 2^{-p}.
\]
\end{definition}

The first form is convenient because a geometric-series calculation gives a tail modulus.  The sequence is never treated as a completed set-theoretic real inside iPRA.  A query asks for \(a(p)\), and the law permits error transport.

\subsection{Arithmetic on names}

Each operation is a demand contract.

\paragraph{Addition and subtraction.}
To output precision \(p\), request both inputs at precision \(p+c\), where \(c\) pays for the triangle inequality and coding convention.  The output term is the rational sum.  The local proof is an integer/rational inequality.

\paragraph{Multiplication.}
Besides input precision, multiplication needs magnitude bounds \(|x|\le B_x\), \(|y|\le B_y\).  Then
\[
|xy-x'y'|\le |x|\,|y-y'|+|y'|\,|x-x'|.
\]
The demand transformer chooses input precisions making the right side smaller than the requested error.  Bounds are part of the presentation or residual producer.

\paragraph{Inversion.}
A represented real is not invertible merely from \(x\ne0\) stated negatively.  The constructive contract requires an apartness witness
\[
  |x|\ge 2^{-k}
\]
or another positive rational lower bound.  If \(|x-x'|\le 2^{-k-1}\), then \(|x'|\ge2^{-k-1}\), and
\[
\left|\frac1x-\frac1{x'}\right|
=\frac{|x-x'|}{|xx'|}
\]
has an explicit bound.  The demand transformer depends monotonically on \(k\) and the output precision.

\paragraph{Composition.}
To compute \(f(g(x))\), use a modulus of continuity for \(f\) on a bounded range containing the approximations to \(g(x)\).  Reverse demand first asks what input accuracy \(f\) needs, then asks \(g\) for that accuracy.  Cyclic dependencies among range bounds and precision are handled by capped demand closure.

\subsection{Apartness and stable predicates}

Exact equality of real names is generally not decidable.  SCOPE therefore distinguishes:

\begin{itemize}
\item metric statements \(|x-y|\le\varepsilon\);
\item apartness statements \(|x-y|\ge\delta>0\);
\item exact algebraic identities inherited from common defining data;
\item discrete labels stabilized by a gap.
\end{itemize}

A sign theorem needs \(|x|\ge\delta\) and an approximation error below \(\delta\).  A branch theorem needs separation.  An exact equality may be proved by an algebraic certificate even when arbitrary equality of represented reals is undecidable.

\subsection{From finite clouds to coherent names}

Non-unique finite approximants do not automatically form one sequence.  A sufficient constructive interface is the following.

For each precision \(p\), let \(C_p\) be a nonempty finite cloud of valid rational codes.  Let \(E_p(a,b)\) be a compatibility relation from \(C_p\) to \(C_{p+1}\).  Assume:

\begin{enumerate}
\item a bounded valid element of \(C_0\) exists;
\item every valid \(a\in C_p\) has a bounded valid \(E_p\)-successor in \(C_{p+1}\);
\item \(E_p(a,b)\) implies \(|a-b|\le m(p)\);
\item a primitive-recursive tail bound satisfies \(\sum_{j\ge p}m(j)\le 2^{-p}\) in the chosen rational coding.
\end{enumerate}

\begin{theorem}[Primitive-recursive coherent selector]
Under these assumptions, the recursion that chooses the least valid initial code and then the least valid compatible successor is primitive recursive and forms a regular Cauchy name.
\end{theorem}
\begin{proof}
Bounded least search supplies the initial code.  Given the current code, the total-extension assumption and bounded search supply a successor.  Primitive recursion therefore defines the sequence.  For \(m>n\), repeatedly apply the edge motion bounds and the triangle inequality:
\[
  |a(m)-a(n)|\le\sum_{j=n}^{m-1}m(j)\le\sum_{j\ge n}m(j)\le2^{-n}.
\]
All sums are finite at each query, and the tail estimate is an explicit object-theory inequality.
\end{proof}

If only ``every finite horizon has some path'' is known, this theorem does not apply.  One must either remain at finite-horizon relational outputs, prove a total extension rule, or supply a shrinking-diameter argument strong enough to make independent finite choices Cauchy.  SCOPE never smuggles an infinite branch out of finite path existence.

\subsection{Metastability versus a Cauchy modulus}

A metastable theorem has the form
\[
  \forall p\,\forall g\,\exists N\le\Phi(p,g)\;
  \forall m,n\in[N,N+g(N)]\ |a_m-a_n|\le2^{-p}.
\]
This is a finite-window theorem and is naturally expressible by bounded reflection and strategy certificates.  It is not a tail modulus
\[
  \forall p\,\exists M\,\forall m,n\ge M\ |a_m-a_n|\le2^{-p}.
\]
No generic SCOPE rule converts the first to the second.  Applications must either use metastability directly or provide additional structure yielding a genuine modulus.

\subsection{Finite compactness interfaces}

Constructive uses of compactness should expose finite data.

\begin{definition}[Net contract]
At mesh \(\eta\), a net contract provides a finite list \(x_0,\ldots,x_{N-1}\) and a coverage verifier proving every represented input is within \(\eta\) of some \(x_i\).  A modulus of uniform continuity transfers a property checked on the net to the full domain with an explicit error.
\end{definition}

The coverage producer may be straightforward for a rational box, or a substantial residual for a complicated space.  Once the net and modulus are supplied, the lifting proof is finite and generic.  The same pattern covers finite subcovers, partitions, triangulations, and quadrature meshes.

\subsection{Finite sums, integrals, and limits}

A Riemann or quadrature argument has three layers:

\begin{enumerate}
\item a finite sum identity or estimate, handled by checked bounded sums and arithmetic;
\item a mesh-error theorem with an explicit modulus;
\item a query compiler choosing mesh size from output precision.
\end{enumerate}

Improper integrals and infinite series need a tail producer.  The finite partial sum is generic; the theorem-specific work is the primitive-recursive tail bound.  SCOPE records it as a residual until proved.

\subsection{Height bounds and complete finite lists}

Let \(P(x)\) be a decidable predicate on natural codes.

\begin{theorem}[Height bound--list equivalence]
\label{thm:height-list}
The following forms of effective finiteness are constructively interconvertible:

\begin{enumerate}
\item an explicit bound \(H\) with \(P(x)\to x<H\);
\item a complete finite list \(L\) such that \(P(x)\) iff \(x\) occurs in \(L\).
\end{enumerate}
\end{theorem}
\begin{proof}
From a bound \(H\), enumerate \(0,\ldots,H-1\) and retain exactly the codes with reflected bit \(P=1\).  Every retained code satisfies \(P\), and every \(P\)-code lies below \(H\) and is retained.  Filtering a bounded list is primitive recursive.

From a complete list, let \(H=0\) for the empty list and otherwise one plus the maximum entry.  Every list member lies below \(H\).  Completeness then gives \(P(x)\to x<H\).  Maximum and membership are primitive recursive, and the proof is finite induction on the list code.
\end{proof}

A bound on the \emph{number} of solutions does not imply either form.  The single solution could have arbitrarily large code.  This is why a cardinality theorem cannot be used as a search cutoff without an additional height argument.

\subsection{Conditional enumerations}

For a Diophantine equation, algebraic variety, finite-state system, or proof search, SCOPE commonly produces
\[
  \mathsf{HeightBound}(H),\Gamma\vdash_{\PRAAx}
  \mathsf{CompleteList}(\mathsf{filter}_{x<H}P(x)).
\]
This separates the deep producer from the routine enumeration and verification.  A later bound theorem closes the residual by simultaneous cut.

\section{Domain-specific embedding blueprints}

The same SCOPE kernel supports different representations.  This section records what the generic framework supplies and what remains domain-specific.

\subsection{Finite algebra and polynomial computation}

Polynomial identities over integers or rationals reduce to finite coefficient equations.  The existing quotient-ring and tactic support can prove many such identities directly in the object theory.

Factorizations, resultants, determinants, and finite Groebner-style reductions are natural certificate nodes.  A certificate contains coefficient arrays and reduction multipliers; the verifier recomputes the identity.  A Nullstellensatz certificate proves infeasibility by checking a finite polynomial combination equal to one.

Roots require care.  Stable public outputs include:

\begin{itemize}
\item isolating intervals with sign-change or derivative certificates;
\item unordered lists or multisets of intervals;
\item symmetric functions of roots;
\item a labeled root only under a separation gap.
\end{itemize}

A complete list of rational or integral solutions additionally needs a height bound.  The finite enumeration thereafter is an instance of \cref{thm:height-list}.

\subsection{Analysis and special functions}

A special function should be presented by one or more certified finite algorithms: a convergent series with tail bound, a continued fraction with error estimate, a differential equation integrator with local truncation bound, or a monotone root bracket.

The SCOPE query records the point, branch region, and output precision.  Reverse demand selects input precision, truncation length, mesh, and internal arithmetic bounds.  Branch-sensitive algorithms carry explicit domain and gap guards.  Overlapping algorithms are combined by a finite cover plus stable-observer agreement on overlaps.

Derivative and integral statements are quantitative.  For example, a differentiability contract supplies a modulus bounding the remainder
\[
  |f(x+h)-f(x)-h f'(x)|\le |h|\,\omega(|h|),
\]
with a primitive-recursive \(\omega\).  An integral contract supplies finite sums and a mesh modulus.  Qualitative phrases such as ``sufficiently small'' or ``uniformly on compacta'' must be replaced by explicit producers.

\subsection{Geometry and tilings}

Finite Euclidean geometry can use integer or rational coordinates, determinant signs, squared distances, and incidence bits.  Degenerate cases are handled relationally rather than by unstable orientation choices.

For tilings and local configurations, a finite patch is a bounded code.  Local matching, coverage of a requested ball, and finite-torus periodicity are decidable verifiers.  Aperiodicity is naturally expressed as rejection of every bounded period code, one finite query at a time.  Coherent patch growth requires compatibility or substitution certificates; a global infinite tiling object need not be internalized if the theorem asks for arbitrarily large patches.

Finite complexes and meshes fit the same pattern.  Homology matrices, incidence relations, and chain certificates are finite algebra.  Topological conclusions requiring compactness or local-to-global passage need a finite cover and explicit overlap data.

\subsection{Probability and expectation}

A finite distribution is a rational weight table with nonnegative entries summing to one.  Expectation, variance, and finite moments are bounded sums.  Conditional expectation on a finite partition is a table computation with guards for zero denominators.

For infinite distributions, a query uses a truncation plus a tail-mass bound.  Concentration inequalities become finite algebraic estimates once moment or exponential-moment producers are supplied.  Almost-sure or measure-theoretic statements should be reformulated as quantitative bounds on finite bad-event probabilities unless a separate constructive measure presentation is formalized.

Randomized algorithms can be treated by finite seed spaces at a query.  A strategy certificate may encode a good seed, while expectation or success probability is a finite rational sum.  Derandomization by least good seed is legitimate only with an explicit finite seed bound and a proof that a good seed exists.

\subsection{Elementary and analytic number theory}

Divisibility, congruences, bounded prime searches, sums, products, and finite combinatorial identities are already close to the checked arithmetic core.  The framework mainly contributes modular certificate composition and finite strategy extraction.

Analytic number theory separates finite identities from analytic producers.  A prime-counting estimate at \(x\) may use a finite combinatorial identity plus explicit remainder constants.  Zero-free regions, exceptional-zero exclusions, and effective constants remain residuals until proved in iPRA.  SCOPE prevents a qualitative classical theorem from being cited as though it supplied those constants.

\subsection{Diophantine geometry}

For a fixed height \(H\), rational or integral points are finitely enumerable and their defining equations are decidable.  Descent certificates, local solubility tables, and finite coverings can be verified.  The generic endpoint is:
\[
  \mathsf{CompleteHeightBound}(H),\mathsf{CoverData}(c),\Gamma
  \vdash_{\PRAAx}\mathsf{CompletePointList}(L).
\]

The hard theorem is typically the uniform production of \(H\) or the complete covering data.  SCOPE does not hide this boundary.  Special families with explicit descent or Runge bounds can close it; general Faltings-, Siegel-, or Mordell--Weil-type producers remain separate research programs.

\subsection{Combinatorics and finite structures}

Finite structures are ideal SCOPE objects.  A graph, poset, matroid, finite category, or automaton is a bounded relation table.  Isomorphisms, colorings, matchings, flows, and decompositions have finite certificate verifiers.  Extremal statements can be represented as bounded games or SAT-style unsatisfiability certificates.

Asymptotic combinatorics needs explicit thresholds and error functions.  A proof that an inequality eventually holds is useful to iPRA only when the threshold is primitive recursive or retained as a residual.

\subsection{Finite category and local-to-global methods}

A fixed finite category is encoded by object and morphism bounds, source/target maps, identity codes, and a composition table with associativity certificates.  Functors, natural transformations, finite limits, and colimits become finite table verifiers.

For local-to-global arguments, one should ask what observation is actually needed.  A full sheaf object may be unnecessary.  A finite cover, overlap compatibility, and spanning-tree descent can establish a root observation or an existential global representative.  Non-tree cycle conditions are additional finite certificate fields.  Exact gluing without canonical representatives should remain existential or quotient-valued.

\subsection{Computability and machine verification}

A finite Turing machine, register machine, or proof checker is a fixed transition table.  Bounded execution is a trace-reflection theorem.  Nonhalting equivalences to consistency statements require a separate metamathematical reduction, but the literal machine transition verification is finite.

When compressing a machine or proof checker, SCOPE recommends keeping the logical endpoint in ordinary iPRA and treating the low-level implementation as a certificate-producing transducer.  State-sharing tricks should be accompanied by a local simulation invariant and a checked translation of traces.

\subsection{A domain checklist}

For each application, answer:

\begin{enumerate}
\item What is the finite query?
\item What codes the input and output?
\item Which predicates are reflected bits?
\item Which loops have QF invariants?
\item What are the resource demands and caps?
\item Are witnesses private, existential, stable, canonical, or residual?
\item What gap, modulus, cover, or height data is genuinely needed?
\item What exact \(\Prf\,\PRAAx\) sequent is the endpoint?
\end{enumerate}

A packet that answers these questions has a clear path to complete formalization.

\section{Resolution of recurring difficulties}

The repository's reports identify both proof-engineering gaps and genuine mathematical barriers.  SCOPE resolves the former and makes the latter explicit.

\small
\begin{longtable}{@{}p{0.22\textwidth}p{0.18\textwidth}p{0.50\textwidth}@{}}
\toprule
Difficulty & Status & Resolution or exact boundary \\
\midrule
\endfirsthead
\toprule
Difficulty & Status & Resolution or exact boundary \\
\midrule
\endhead
Semantic Lean theorem mistaken for an iPRA proof & Resolved & Public endpoints are values of \(\Prf\,\PRAAx\).  Every graph node carries an object derivation builder; Python and semantic Lean code are discovery or test layers only. \\

Quantifiers occur in an induction invariant & Resolved & Compile bounded properties to characteristic terms.  Induct on equations about bits, traces, ranks, and accumulators.  Reflect back before or after induction. \\

Nested bounded quantifier alternation & Resolved & Fixed bounded formulas have reflected bits.  When choices are needed, finite strategy and counterstrategy tables are checked by bounded games. \\

An all-precision invariant is demanded in one loop & Resolved by reformulation & Fix the terminal query first, close finite backward demands, and run one forward trace.  Do not quantify over all future precisions inside the induction matrix. \\

Cyclic precision dependencies & Resolved & Iterate the monotone inflationary demand map on an explicitly capped product.  Rank gives a fixed point by total capacity and proves leastness. \\

Least local witness is a dead end & Resolved & Compute backward viability before selecting the least forward continuation. \\

Proof-local certificate leaks into a public theorem & Resolved & Apply the five-mode escape analysis: private, existential, stable, canonical, or residual. \\

Nonunique approximants are treated as one coherent infinite name & Conditional & Keep a relational finite output.  A coherent name requires a total extension rule with a summable motion budget, or another explicit shrinking-diameter construction. \\

Exact ordering of colliding roots or eigenvectors & Barrier with correct replacement & Export an unordered multiset, symmetric invariant, projector, or finite relation.  Exact labels require a separation gap. \\

Division by a merely nonzero represented real & Barrier with correct replacement & Require an apartness witness \(|x|\ge\delta>0\).  Negative nonzeroness alone gives no modulus for inversion. \\

Compactness invoked without finite data & Conditional & Supply a finite net or cover, a coverage verifier, and a modulus transferring local checks to the domain. \\

Metastability is treated as a convergence rate & Barrier & Preserve the finite-window theorem.  A Cauchy name needs a separate tail modulus or stronger structure. \\

Cardinality bound is used as a search cutoff & Barrier & A cardinality bound gives no height.  Require a height cutoff or complete finite list; these are constructively interconvertible for decidable codes. \\

Qualitative uniqueness is treated as computable selection & Barrier & Require quantitative determinacy: isolation, contraction, strict convexity with a modulus, or a finite gap. \\

A classical proof is silently imported & Conditional & Use a finite checkable certificate, or formalize a negative translation for the precise fragment and stable conclusion.  Classical discovery does not itself produce an object proof. \\

Open hypotheses are lost in backend translation & Resolved in the current repository & Use the checked open and scoped \code{IPRABridge} transports.  Develop reusable results in ordinary \(\Prf\,\PRAAx\). \\

Arity-indexed abstraction of open terms & Resolved in the current repository & Reuse \code{PTm.toPR} and \code{toPR_eq}; use \code{recDef} for parameterized folds. \\

A universal evaluator for arbitrary PR codes is requested & Rejected as unnecessary & Specialize every fixed program in Lean's metatheory.  SCOPE is a schema, not an internal universal interpreter. \\

Deep effective Diophantine finiteness & Not generically solved & The finite enumeration/checker is generic once a height producer exists.  The effective height theorem remains the residual. \\

Analytic constants or zero-free regions are missing & Not generically solved & The finite analytic trace accepts explicit constants and moduli.  Producing them is theorem-specific mathematics. \\
\bottomrule
\end{longtable}
\normalsize

\subsection{Classical-to-intuitionistic transfer}

Some classical proofs can be transferred by negative translation, especially when their atoms are decidable and their conclusions are stable.  This is useful but should not be overgeneralized.

For a fixed finite bounded argument, direct Boolean reflection is usually simpler: it supplies positive witnesses, least counterexamples, and executable verifiers.  A negative translation is better suited to an already formalized proof in a compatible classical arithmetic where the translated axioms and induction rules have been checked.  SCOPE treats this as a separate compiler path with its own correctness theorem.

\subsection{Why the barriers are mathematically informative}

The barriers identify the missing modulus or producer that a future proof must deliver.  They are not reasons to abandon the theorem.  For example:

\begin{itemize}
\item replacing root labels by multisets may reveal the invariant actually used downstream;
\item isolating an apartness margin may turn a qualitative proof into a quantitative algorithm;
\item asking for a height cutoff exposes the exact effective content of a finiteness theorem;
\item preserving metastability may be enough for the intended finite computation.
\end{itemize}

A framework is successful when failed compilation produces a mathematically precise missing contract rather than an opaque proof-assistant obstacle.

\section{Worked finite compilations}

The examples in this section are small enough to inspect manually, but each exercises a distinct generic layer.  Their exact host-level computations are recorded in \path{examples/demo_certificate.json}.

\subsection[A bounded forall-exists strategy]{A bounded \(\forall\exists\) strategy}

Consider
\[
  \forall x<6\;\exists y<4\quad (x+y)\bmod 3=0.
\]
The matrix is decidable by the checked division/remainder arithmetic.  For each \(x\), bounded least search returns the least \(y<4\) satisfying the congruence.  The resulting table is
\[
  \sigma=[0,2,1,0,2,1].
\]
Indeed,
\[
0+0,\ 1+2,\ 2+1,\ 3+0,\ 4+2,\ 5+1
\]
are all divisible by three.  The table verifier checks six bounds and six remainder equations.  By finite strategy soundness, the verifier implies the quantified formula.

There are several possible public endpoints:

\begin{itemize}
\item prove the bounded formula and keep \(\sigma\) private;
\item return \(\exists\sigma\,\mathsf{Verifies}(\sigma)\);
\item return the least strategy table by canonical search;
\item export only a stable statistic of the table.
\end{itemize}

The choice among them is a theorem-design decision, not an implementation accident.

\subsection{A cyclic precision network}

Take caps \((5,6,5)\) and demand transformer
\[
F(d_0,d_1,d_2)=
\bigl(\min(5,d_1+1),\min(6,d_2+2),\min(5,d_0)\bigr).
\]
It is monotone.  Starting at zero and applying \(H(d)=d\vee F(d)\) gives
\[
\begin{array}{c|c|c}
 t & d_t & \rank(d_t)\\
\hline
0&(0,0,0)&0\\
1&(1,2,0)&3\\
2&(3,2,1)&6\\
3&(3,3,3)&9\\
4&(4,5,3)&12\\
5&(5,5,4)&14\\
6&(5,6,5)&16\\
7&(5,6,5)&16.
\end{array}
\]
The total cap rank is \(5+6+5=16\), and the fixed point is reached after six strict iterations.  The example demonstrates why one reverse pass is insufficient and why the rank budget gives a uniform stopping bound.

In an analytic graph, the coordinates could represent input precision, range-bound precision, and a denominator-separation precision.  The fixed point is computed before the forward numerical trace, so the trace invariant mentions only concrete resource terms.

\subsection{Viability versus greedy choice}

Let
\[
A_0=\{0,1\},\qquad A_1=\{0,1\},\qquad A_2=\{0\}.
\]
All nodes are locally valid.  The compatibility edges are
\[
  0\to0,\quad 1\to1
  \qquad\text{between }A_0,A_1,
\]
and only
\[
  1\to0
  \qquad\text{between }A_1,A_2.
\]
Greedy least local choice starts with \(0\), moves to \(0\), and fails.  Backward viability yields
\[
  V_2=\{0\},\qquad V_1=\{1\},\qquad V_0=\{1\}.
\]
The least viable path is \((1,1,0)\).  The example is deliberately included in the formalization acceptance gates because it catches a subtle but common incorrect selector.

\subsection{Metric observation on a finite tree}

Take a rooted tree with edges \(r\to a\), \(r\to b\), and \(a\to c\).  Suppose certified edge errors are
\[
 e_a=2,\qquad e_b=1,\qquad e_c=3.
\]
Metric tree descent gives root-distance bounds
\[
  E_r=0,
  \quad E_a=2,
  \quad E_b=1,
  \quad E_c=e_a+e_c=5.
\]
The public result can expose the root observation and these explicit error radii.  It need not choose globally compatible representatives on every cycle of a larger cover unless the downstream theorem actually requires them.

\subsection{Height-bounded Diophantine enumeration}

Let \(P(x)\) be the reflected equation predicate for a fixed Diophantine problem under a natural code.  Assume a residual formula
\[
  \mathsf{HB}(H):=\forall x\,[P(x)\to x<H].
\]
The compiler filters \(0,\ldots,H-1\), producing a list \(L_H\).  Bounded reflection proves:
\[
  x\in L_H\leftrightarrow x<H\wedge P(x).
\]
Under \(\mathsf{HB}(H)\), this becomes \(x\in L_H\leftrightarrow P(x)\).  The result is the conditional object theorem
\[
  \mathsf{HB}(H)\vdash_{\PRAAx}\mathsf{Complete}(L_H,P).
\]

Nothing in the finite compiler claims to produce \(H\).  A theorem-specific descent, Runge estimate, Baker bound, or other effective argument can later discharge the residual.

\subsection{A quantitative bisection contract}

Let \(f\) be a fixed primitive-recursive rational function on a rational interval \([a,b]\).  Suppose the residual context supplies:

\begin{enumerate}
\item opposite endpoint signs;
\item a monotonicity certificate on the interval;
\item a denominator/apartness guard for every rational operation;
\item a primitive-recursive bound \(N(p)\) such that \((b-a)/2^{N(p)}\le2^{-p}\).
\end{enumerate}

A trace stores endpoints \((\ell_t,r_t)\).  At each step it computes the midpoint and uses the reflected sign bit to retain a bracketing half.  The quantifier-free invariant states:
\[
 a\le\ell_t\le r_t\le b,
 \quad f(\ell_t)\le0\le f(r_t),
 \quad 2^t(r_t-\ell_t)=b-a.
\]
After \(N(p)\) steps, return the midpoint and radius \(2^{-p}\).  Monotonicity proves every exact root in the bracket lies in the output ball; a separation/strict-monotonicity modulus can strengthen this to uniqueness.  The theorem constructs a certified finite approximation without postulating a completed real root.

\subsection{End-to-end composition}

A representative SCOPE compilation has the following order:

\begin{enumerate}
\item specialize the theorem graph to query \(q\);
\item compute capped cyclic demand closure;
\item decode or produce a bounded master certificate;
\item reflect all bounded validity conditions;
\item compute viability for nonunique paths;
\item run fixed forward traces at the closed resources;
\item project through an exact or metric observer;
\item eliminate private certificates;
\item retain only reachable producer residuals;
\item emit the final \(\Prf\,\PRAAx\) derivation.
\end{enumerate}

The demo JSON records finite evidence for steps 2, 5, 6 in strategy form, and 7 in metric-tree form.  It is executable evidence for the algorithms, not a substitute for the object derivations described above.

\section{Lean formalization architecture}
\label{sec:lean-plan}

\subsection{Integration target}

The ordinary implementation should live under
\begin{center}
\code{Logic/PrimitiveRecursiveArithmetic/Derivations/Lean/}\allowbreak
\code{IPRADerivations/Framework/SCOPE/}.
\end{center}
It should be exposed by an umbrella import only after every module builds.  Foundation adapters belong in \code{IPRABridge} and should reuse the already checked open-sequent transports.

\subsection{Structural kernel}

The archive's \code{lean/SCOPE/StructuralKernel.lean} is a no-placeholder source-level candidate.  Its core declarations have the following shapes:

\begin{lstlisting}
theorem replaceContext (h : Gamma |- A)
    (supply : forall C in Gamma, Delta |- C) : Delta |- A :=
  h.trans supply

structure Contract (n) (Gamma) (goal) where
  valid   : Fm PTm (n + 1)
  extract : valid :: Gamma.map Fm.shift |- goal.shift

theorem Contract.close
    (producer : Gamma |- Fm.ex C.valid) : Gamma |- goal :=
  .exE producer C.extract

structure ResidualResult (n) (Gamma) (goal) where
  residuals : List (Fm PTm n)
  derivation : residuals ++ Gamma |- goal
\end{lstlisting}

The actual file uses the repository notation \code{Gamma |- [PRAAx] A}.  It adds residual closure, weakening, specialization, boundary introduction, and direct sealing.  It contains no \code{sorry} or new axiom.  It was not elaborated in this runtime because the repository workspace was unavailable locally, so it is not reported as checked.

\subsection{Bounded formula module}

The first substantial module should define the fixed syntax:

\begin{lstlisting}
inductive BForm : Nat -> Type
  | eq    (s t : PTm n)
  | bot
  | and   (A B : BForm n)
  | or    (A B : BForm n)
  | imp   (A B : BForm n)
  | allLt (bound : PTm n) (body : BForm (n+1))
  | exLt  (bound : PTm n) (body : BForm (n+1))
\end{lstlisting}

Define \code{denote} in \(\Fm\) and \code{bit} in \(\PTm\).  The key arbitrary-arity fold construction is:

\begin{lstlisting}
def liftFoldBody (body : PTm (n+1)) : PTm (n+2) :=
  Syn.subst (Fin.cases (var 0) (fun j => var j.succ.succ)) body

def allFoldPR (body : PTm (n+1)) : PR (n+1) :=
  recDef (num 1) (band (var 1) (liftFoldBody body))
\end{lstlisting}

This explicitly resolves the parameter-arity problem.  The compiler should return a package:

\begin{lstlisting}
structure Reflection (A : Fm PTm n) where
  bit      : PTm n
  boolean  : [] |- (bit = 0 or bit = 1)
  forward  : [] |- (A -> bit = 1)
  backward : [] |- (bit = 1 -> A)
\end{lstlisting}

The implementation order is Boolean algebra, fold equations, search adapters, fold reflection, then structural recursion on \code{BForm}.  Nested binder tests must substitute non-variable terms through at least two alternating quantifiers.

\subsection{Least accepted adapter}

Wrap the existing search API in the accepted-bit convention.  A useful structure is:

\begin{lstlisting}
structure LeastAccepted (accepted : PTm (n+1)) (bound : PTm n) where
  witness  : PTm n
  foundBit : PTm n
  boolean  : ...
  sound    : ...
  minimal  : ...
  sentinel : ...
\end{lstlisting}

The implementation should reuse \code{Search.search_found}, \code{search_minimal}, and \code{search_spec}; it should not reprove least search from scratch.

\subsection{Codes and demand closure}

Select one vector representation and prove it before implementing closure.  Mixed radix is attractive because total capacity and lexicographic code order are explicit.  Required modules are:

\begin{enumerate}
\item \code{MixedRadix}: digit, encode, decode, range, and round trip;
\item \code{CappedVector}: coordinate order, join, rank, and cap sum;
\item \code{DemandClosure}: inflationary iteration, strict rank, fixedness, and leastness.
\end{enumerate}

The finite dimension is a Lean parameter or a fixed structure index; object iteration is on the natural step counter.  Coordinatewise boundedness is represented by a reflected fold bit so the induction formula remains QF.

\subsection{Trace module}

A reusable trace package should accept fixed terms/symbols and local derivations rather than a semantic Lean function alone.  One possible interface is:

\begin{lstlisting}
structure TraceContract (n : Nat) where
  statePR  : PR (n+1)
  init     : PTm n
  inv      : Fm PTm (n+1)
  invQF    : inv.QF
  base     : Gamma |- inv[init]
  step     : inv :: Gamma.shift |- inv[next]
\end{lstlisting}

The precise indexing should be chosen to minimize rewrites with \code{recDef}.  The result theorem applies \code{IPRA.ind} and then \code{Prf.subst} at arbitrary fuel.

\subsection{Finite strategies and viability}

Begin with the \(\forall\exists\) case.  Encode a table by mixed radix and prove table lookup equations.  Then add a host-level datatype for fixed quantifier shapes and compile it recursively.

Viability should be a separate module.  Its acceptance test must include the dead-end example from \cref{sec:viability}.  Soundness and completeness should be stated as ordinary object derivations of reflected path predicates, not merely as Lean propositions over lists.

\subsection{Observers and scope effects}

Define structures for:

\begin{itemize}
\item exact observer to a fixed outer term;
\item explicit congruence observer;
\item rational metric diameter observer;
\item gap-guarded finite label observer.
\end{itemize}

The scope compiler can use a host enum
\begin{lstlisting}
inductive Escape
  | private | existential | stable | canonical | residual
\end{lstlisting}
but each constructor must require the proof payload corresponding to its semantics.  In particular, \code{stable} is not a Boolean annotation; it carries an observer derivation.

\subsection{Graph compiler}

The graph compiler is a Lean metaprogram or ordinary recursive function over a fixed typed graph.  Every node returns a term, formula, and derivation package.  The final result is a \code{ResidualResult}.  The compiler may fail with a structured diagnostic when a resource cap, observer proof, or producer is missing.

A useful diagnostic is the minimal residual frontier: traverse backward from the goal, remove producers discharged internally, and report only reachable unsupplied leaves.  This makes failed compilation mathematically informative.

\subsection{Acceptance gates}

The formalization plan in the archive lists the complete gates.  The most important are:

\begin{itemize}
\item focused and full \code{lake build} success;
\item no \code{sorry}, \code{admit}, new \code{axiom}, or unsafe Boolean reflection;
\item \code{#print axioms} audit consistent with repository policy;
\item substitution at compound terms;
\item nested bounded alternation;
\item a cyclic demand example requiring multiple passes;
\item a greedy-dead-end viability example;
\item two distinct valid certificates with one stable public observation;
\item an exact residual visible in the final theorem type;
\item replay through both scoped bridge directions.
\end{itemize}

\subsection{Four staged pilots}

After the core builds, migrate:

\begin{enumerate}
\item a finite arithmetic theorem with a bounded certificate;
\item the quantitative bisection contract of \cref{sec:trace};
\item a finite complex or cover example using tree descent;
\item a height-bounded Diophantine list with the height theorem left residual.
\end{enumerate}

These pilots exercise every layer without requiring an immediate attack on the hardest analytic or Diophantine producers.

\section{Sharp barriers and nonclaims}

A reliable framework must state not only what it compiles, but why certain stronger requests cannot be obtained from the advertised data.

\subsection{No universal primitive-recursive evaluator}

Suppose there were a primitive-recursive total function \(U(e,x)\) evaluating every unary primitive-recursive program code \(e\).  Then the diagonal function
\[
  d(e)=U(e,e)+1
\]
would be primitive recursive and therefore have some code \(e_d\).  But
\[
  d(e_d)=U(e_d,e_d)+1=d(e_d)+1,
\]
a contradiction.  SCOPE avoids this diagonal obstruction by compiling each fixed term/program in Lean's metatheory.

\subsection{Discontinuous exact selectors}

Ordering roots through a collision, choosing a sign for a zero eigenvector, or deciding the sign of an arbitrarily small represented real are discontinuous operations.  Finite approximations without a gap cannot determine them uniformly.  A proof assistant cannot repair missing information.  The correct outputs are symmetric, quotient, relational, or metric until a separation witness is supplied.

\subsection{Nonzero versus apart}

In intuitionistic analysis, \(x\ne0\) means that \(x=0\) leads to contradiction.  It does not necessarily provide a rational \(\delta>0\) with \(|x|\ge\delta\).  Computing \(1/x\) to a requested precision requires such a lower bound because the conditioning worsens near zero.  Therefore inversion contracts explicitly require apartness.

\subsection{Cardinality does not give height}

For every proposed cutoff function \(H(N)\), the singleton predicate \(P(x)\equiv x=H(1)+1\) has cardinality one but its solution lies beyond \(H(1)\).  Thus even a perfect cardinality bound does not imply an enumeration bound.  A height theorem or complete list is strictly stronger effective data.

\subsection{Metastability does not give a modulus}

A metastable bound controls each finite window chosen by a challenge function.  Turning these independently located windows into one tail index requires additional uniformity.  SCOPE does not assert such a conversion.  This prevents a common but invalid promotion of finite convergence data to a regular real name.

\subsection{Qualitative compactness is not a finite net}

A classical compactness theorem may assert that every cover has a finite subcover without supplying a primitive-recursive search space, coverage decision procedure, or modulus.  SCOPE consumes a coded finite cover/net and its verification.  Producing it uniformly is separate constructive content.

\subsection{Finite path existence is not an infinite branch}

The fact that every finite horizon has a compatible path does not by itself supply one primitive-recursive infinite path.  Even classically, extracting a branch uses a compactness principle; constructively it needs a choice/extension mechanism.  SCOPE remains at finite horizons unless a total successor rule, coherent strategy, or shrinking quantitative interface is proved.

\subsection{External checking is not object proof}

A Python verifier can establish that a finite file passes a test in the ambient runtime.  It cannot by itself create a theorem of the embedded arithmetic.  The final trusted path must represent the verifier as primitive-recursive terms and prove its soundness in \(\Prf\,\PRAAx\), or translate a separately checked proof object through a verified bridge.

\subsection{Generic verification does not imply generic production}

A small verifier may check a very deep certificate.  This does not make the certificate easy to produce, nor does it prove that one always exists.  SCOPE's residual formulas preserve this distinction.  Examples include effective Mordell bounds, zero-free regions, complete descent data, and finite covers with quantitative constants.

\subsection{Formal status of this archive}

The present report supplies human proofs and a formalization plan.  The Python layer has been executed and tested.  The structural Lean file was source-audited against current APIs but not elaborated in a local ProveIt workspace.  No new kernel-checked theorem is claimed.  The archive's purpose is to provide a mathematically complete design whose remaining formal work is explicit and dependency ordered.

\section{Further research questions}

\subsection{Immediate formalization questions}

\begin{enumerate}
\item Can the full fixed bounded-formula reflection compiler be implemented with proof terms small enough for routine use, while sharing fold and Boolean lemmas rather than expanding them?
\item Which vector coding---mixed radix, nested pairing, or a small fixed tuple syntax---gives the shortest demand-closure proofs in the current library?
\item Can the graph compiler infer the minimal residual frontier automatically and produce a human-readable explanation of every remaining producer?
\item Can observer stability obligations be synthesized from local equality/metric lemmas along a supplied dependency tree?
\item What is the best typed interface for a node whose output is relational rather than functional?
\item Can generated proof DAGs be hash-consed or reflected through a small checked certificate calculus without changing the ordinary \(\Prf\,\PRAAx\) endpoint?
\end{enumerate}

\subsection{Mathematical framework questions}

\begin{enumerate}
\item Which classical theorems admit a direct SCOPE certificate with primitive-recursive bounds even when their customary proof uses compactness or choice?
\item Can a general ``shrinking cloud'' theorem characterize exactly when finite relational approximants determine a regular Cauchy name?
\item Can metric observer composition be organized as an error semiring so that reverse demand and forward error proofs are generated from one algebraic interface?
\item How far can finite category, diagram, and Cech-style arguments be developed using only explicit covers, compatibility tables, and observational descent?
\item Which root/eigenvalue theorems are most naturally stated using symmetric observations rather than labeled branches, and how much proof complexity is saved?
\item Can proof mining systematically turn metastable packets into SCOPE finite strategies with explicit residual moduli?
\end{enumerate}

\subsection{Producer-focused questions}

\begin{enumerate}
\item For which Diophantine families can explicit primitive-recursive height bounds be proved in iPRA using the existing descent and divisibility libraries?
\item Can effective Runge, Thue, Pell, or special elliptic-curve bounds be isolated as reusable producer modules?
\item Which analytic number-theory estimates can be reduced to finite interval arithmetic plus a small list of explicit residual constants?
\item Can finite-net producers for common compact domains be standardized with checked coverage and modulus interfaces?
\item Can apartness and branch-gap producers be extracted automatically from interval or derivative certificates?
\end{enumerate}

\subsection{Logic and translation questions}

\begin{enumerate}
\item What is the largest useful classical fragment whose negative translation can be implemented and verified once, while preserving the finite witness information needed by SCOPE?
\item Can the ordinary/native/ Foundation translations preserve SCOPE residual and observer structures definitionally, or only propositionally?
\item Is there a concise internal calculus specialized to certificate sealing and reflected finite traces whose proofs translate linearly to ordinary \(\Prf\,\PRAAx\)?
\item Can the current concise native calculi reduce proof-checker size for generated SCOPE certificates without complicating application development?
\end{enumerate}

\subsection{Suggested first research sequence}

A high-leverage sequence is:

\begin{enumerate}
\item check the structural kernel and bounded Boolean algebra;
\item complete arbitrary-arity fold reflection;
\item implement least accepted search and one \(\forall\exists\) strategy;
\item implement mixed-radix demand closure;
\item formalize the dead-end viability example;
\item add exact and metric observer elimination;
\item build the four pilots from \cref{sec:lean-plan};
\item only then generalize the host graph compiler.
\end{enumerate}

This order produces useful checked mathematics at every milestone and avoids designing an elaborate DSL before the core object lemmas are stable.

\appendix
\section{Consolidated theorem interfaces}

This appendix collects the central interfaces in a form suitable for a formalization ledger.

\subsection{Structural}

\begin{enumerate}
\item \textbf{Residual cut:}
\[
R,\Gamma\vdash A,
\quad (\forall B\in R\;\Gamma\vdash B)
\quad\Longrightarrow\quad
\Gamma\vdash A.
\]
\item \textbf{Private seal:}
\[
\Gamma\vdash\exists c\,V(c),
\quad V(c),\Gamma^+\vdash A^+
\quad\Longrightarrow\quad
\Gamma\vdash A.
\]
\item \textbf{Specialization:}
\[
\Gamma(\vec x)\vdash A(\vec x)
\quad\Longrightarrow\quad
\Gamma(\vec t)\vdash A(\vec t).
\]
\end{enumerate}

\subsection{Bounded reflection}

For every fixed bounded formula \(A\):
\[
\vdash \bit_A=0\vee\bit_A=1,
\qquad
\vdash A\leftrightarrow\bit_A=1.
\]
For a bounded relation \(R(x,y)\), least accepted search provides \(w,h\) with:
\[
 h=1\leftrightarrow\exists y<B\,R(x,y),
\quad h=1\to w<B\wedge R(x,w),
\quad h=0\to w=B,
\]
and leastness below \(w\).

\subsection{Trace reflection}

For a fixed primitive-recursive trace and QF invariant:
\[
I(s_0,q),
\quad I(s,q)\to I(T(s,q),q)
\quad\Longrightarrow\quad
I(\mathsf{run}(N,q),q).
\]
All displayed implications stand for object derivations under the appropriate contexts.

\subsection{Demand closure}

For monotone \(F\) on \(D=\prod_i[0,C_i]\), with \(H(d)=d\vee F(d)\) and \(K=\sum_iC_i\):
\[
H^{K+1}(0)=H^K(0),
\qquad
F(e)\le e\to H^K(0)\le e.
\]

\subsection{Viability}

For finite levels and backward bits \(V_i\):
\[
V_i(a)=1
\quad\leftrightarrow\quad
\text{there exists a valid compatible suffix beginning at }a.
\]
A path exists iff some initial viability bit is one, and bounded least search yields the lexicographically least viable path.

\subsection{Observers}

\begin{enumerate}
\item Exact: \(V(c)\to O(c)=o\) lets any invariant consumer export \(o\).
\item Quotient: pairwise congruence lets congruence-respecting consumers eliminate the certificate.
\item Metric: pairwise diameter \(\le\varepsilon\) yields an existential or canonical center with radius \(\varepsilon\).
\item Gap: approximation error \(e\) and exact gap \(g>2e\) stabilize the finite label.
\end{enumerate}

\subsection{Scope normal form}

A fixed finite graph satisfying local contracts compiles to
\[
  B_1(q),\ldots,B_r(q),\Gamma(q)\vdash_{\PRAAx}\Goal(q),
\]
with every internal certificate private, existential, stable, canonical, or residual.  No other escape is permitted.

\section{Repository baseline and artifact inventory}

\subsection{Pinned source snapshot}

The report was prepared from the public ProveIt \code{main} commit
\[
\code{58627ed89814ba122e1f0f012f1aba103cda267b}
\]
dated 16 November 2026 at 20:33:53 UTC.

The principal inspected files were:

\begin{itemize}
\item \code{ConciseIPRA/Lean/ConciseIPRA/Syntax.lean};
\item \code{ConciseIPRA/Lean/ConciseIPRA/Deduction.lean};
\item \code{ConciseIPRA/Lean/ConciseIPRA/PR.lean};
\item \code{Derivations/Lean/IPRADerivations/Core/Basic.lean};
\item \code{Core/Definitions.lean}, \code{Core/Pairing.lean}, and \code{Core/Search.lean};
\item \code{Derivations/README.md} and \code{UNIFICATION-PLAN.md};
\item \code{Bridge/README.md};
\item the foundations catalog in \code{docs/proofs/README.md};
\item the bounded-reflection and scoped-certificate packets of 16 November 2026;
\item earlier unified, PFIG, KIFC-R, and FINER framework packets.
\end{itemize}

The complete path-level audit is in \code{metadata/source_audit.md}.

\subsection{Archive inventory}

\begin{center}
\begin{tabularx}{\textwidth}{@{}lX@{}}
\toprule
Path & Role \\
\midrule
\code{scope_ipra_framework.tex} & Main editable report source. \\
\code{sections/*.tex} & Section sources. \\
\code{scope_ipra_framework.pdf} & Compiled report. \\
\code{lean/SCOPE/StructuralKernel.lean} & No-placeholder structural candidate. \\
\code{lean/SCOPE/*Blueprint.lean} & Main definitions and theorem targets. \\
\code{lean/FORMALIZATION_PLAN.md} & Dependency-ordered implementation plan and gates. \\
\code{reference/scope_reference.py} & Exact finite reference algorithms. \\
\code{reference/test_scope_reference.py} & Fourteen unit tests. \\
\code{examples/demo_certificate.json} & Deterministic worked evidence. \\
\code{metadata/theorem_ledger.csv} & Status of every principal theorem. \\
\code{metadata/difficulty_matrix.csv} & Resolution/barrier classification. \\
\code{metadata/claims.md} & Formal-status and nonclaim policy. \\
\code{metadata/integration_map.md} & Mapping to current ProveIt modules. \\
\code{validation/} & Captured test, preflight, and render outputs. \\
\bottomrule
\end{tabularx}
\end{center}

\subsection{Reproduction}

Run \code{bash build.sh}.  It executes the Python test suite, regenerates the demo JSON, compiles the PDF, runs PDF preflight, and renders every page for visual inspection.  Lean integration requires a ProveIt checkout and follows \code{lean/FORMALIZATION_PLAN.md}.

\subsection{Final status statement}

The archive supplies a complete human mathematical framework, exact finite executable evidence, and a clear path to checked implementation.  It deliberately does not conflate those layers.  Existing repository theorems are inherited as checked; all new generic results remain human proofs and Lean targets until elaborated in the ProveIt workspace.


\end{document}
